% =====================================================================
%  Conclusion : importée par main.tex avec % =====================================================================
%  Conclusion : importée par main.tex avec % =====================================================================
%  Conclusion : importée par main.tex avec % =====================================================================
%  Conclusion : importée par main.tex avec \include{conclusion}.
%  Section non numérotée, ajoutée au sommaire.
%  Sources : les chapitres du cours, le registre des projets officiels
%  (gouvernance d'OpenStack), l'annonce de l'arrivée de l'OpenInfra Foundation
%  à la Linux Foundation, la page de l'examen COA et le guide du contributeur
%  (voir sources.bib). Conclusion détaillée : 10 pages au plus.
% =====================================================================

\phantomsection
\section*{Conclusion}
\addcontentsline{toc}{section}{Conclusion}

% Figures et tableaux de cette section non numérotée : numérotation propre (C.1, C.2\dots)
% pour ne pas prolonger celle du chapitre précédent.
\setcounter{figure}{0}\setcounter{table}{0}
\renewcommand{\thefigure}{C.\arabic{figure}}\renewcommand{\thetable}{C.\arabic{table}}
\renewcommand{\theHfigure}{C.\arabic{figure}}\renewcommand{\theHtable}{C.\arabic{table}}

Ce cours a présenté OpenStack comme une famille de services indépendants qui, ensemble, transforment des serveurs, des disques et
des réseaux en un cloud en libre-service. Il est parti de la carte d'ensemble et de l'identité (Keystone, partie~\ref{part:fondations}),
a suivi la vie d'une machine virtuelle à travers le socle IaaS (Nova, Placement, Glance, Neutron, Cinder, Swift,
partie~\ref{part:socle}), puis a présenté les services qui permettent d'exploiter le cloud (Horizon, Octavia, Heat,
partie~\ref{part:exploiter}) et ceux qui fournissent d'autres ressources au-dessus du socle (Ironic, Magnum, Trove,
partie~\ref{part:plateformes}). Il s'est achevé du côté de l'administrateur, avec le déploiement (partie~\ref{part:miseenoeuvre}). Cette conclusion rassemble l'ensemble : une carte,
un scénario de bout en bout, les principes qui reviennent d'un service à l'autre, des repères pour choisir et exploiter, puis les
limites de la plateforme, ce que le cours n'a pas traité et la manière de poursuivre.

% ---------------------------------------------------------------------
\subsection*{La carte d'ensemble}
% ---------------------------------------------------------------------

La figure~\ref{fig:concl-carte} reprend la carte du chapitre~\ref{chap:presentation} (figure~\ref{fig:pres-archi}) en y ajoutant les
numéros de chapitre. Tous les services reposent sur la même infrastructure (base de données, file de
messages, cache) et sur Keystone, qui authentifie chaque appel ; l'outil de déploiement installe et met à jour le tout. Le
tableau~\ref{tab:concl-services} rappelle, pour chaque service, son rôle et les processus à connaître.

\begin{figure}[!ht]
\centering
\adjustbox{max width=\linewidth}{%
\begin{tikzpicture}[>=Stealth, font=\footnotesize]
  \node[pcli, minimum width=13.0cm] at (8.5,8.3)
    {Utilisateurs \ $\cdot$ \ ligne de commande \ $\cdot$ \ API REST \ $\cdot$ \ Horizon (chap.~\ref{chap:horizon})};
  % services d'exploitation et de plateforme
  \node[psvc, minimum width=2.6cm, minimum height=1.1cm, inner xsep=2pt] at (2.7,6.8) {\textbf{Heat} {\scriptsize\mdseries(chap.~\ref{chap:heat})}\\ {\scriptsize\mdseries orchestration}};
  \node[psvc, minimum width=2.6cm, minimum height=1.1cm, inner xsep=2pt] at (5.6,6.8) {\textbf{Octavia} {\scriptsize\mdseries(chap.~\ref{chap:octavia})}\\ {\scriptsize\mdseries répartition de charge}};
  \node[psvc, minimum width=2.6cm, minimum height=1.1cm, inner xsep=2pt] at (8.5,6.8) {\textbf{Ironic} {\scriptsize\mdseries(chap.~\ref{chap:ironic})}\\ {\scriptsize\mdseries machines physiques}};
  \node[psvc, minimum width=2.6cm, minimum height=1.1cm, inner xsep=2pt] at (11.4,6.8) {\textbf{Magnum} {\scriptsize\mdseries(chap.~\ref{chap:magnum})}\\ {\scriptsize\mdseries Kubernetes managé}};
  \node[psvc, minimum width=2.6cm, minimum height=1.1cm, inner xsep=2pt] at (14.3,6.8) {\textbf{Trove} {\scriptsize\mdseries(chap.~\ref{chap:trove})}\\ {\scriptsize\mdseries bases de données}};
  % services de base
  \node[pfoc, minimum width=3.75cm, minimum height=1.1cm] at (4.5,5.2) {Nova, Placement\\ {\scriptsize\mdseries calcul (chap.~\ref{chap:nova}, \ref{chap:placement})}};
  \node[pfoc, minimum width=3.75cm, minimum height=1.1cm] at (8.5,5.2) {Neutron\\ {\scriptsize\mdseries réseau (chap.~\ref{chap:neutron})}};
  \node[pfoc, minimum width=3.75cm, minimum height=1.1cm] at (12.5,5.2) {Glance\\ {\scriptsize\mdseries images (chap.~\ref{chap:glance})}};
  \node[psvc, minimum width=3.75cm, minimum height=1.1cm] at (6.5,3.75) {Cinder\\ {\scriptsize\mdseries volumes (chap.~\ref{chap:cinder})}};
  \node[psvc, minimum width=3.75cm, minimum height=1.1cm] at (10.5,3.75) {Swift\\ {\scriptsize\mdseries objets (chap.~\ref{chap:swift})}};
  % infrastructure
  \node[pgris, minimum width=13.0cm, minimum height=0.95cm] at (8.5,2.3)
    {Infrastructure partagée\\ {\scriptsize\mdseries MariaDB \ $\cdot$ \ RabbitMQ \ $\cdot$ \ memcached \ $\cdot$ \ HAProxy}};
  \node[pgris, minimum width=13.0cm, minimum height=0.95cm] at (8.5,1.05)
    {Machines et réseaux\\ {\scriptsize\mdseries hyperviseurs, disques, réseau physique, serveurs nus}};
  % barres latérales
  \node[draw=insarouge, thick, rounded corners=4pt, fill=insarouge!8, minimum width=1.5cm, minimum height=8.7cm] at (0.45,4.65) {};
  \node[rotate=90, font=\footnotesize\bfseries, text=insarouge] at (0.45,4.65) {Keystone : identité, catalogue, tokens (chap.~\ref{chap:keystone})};
  \node[draw=insarouge, thick, rounded corners=4pt, fill=insarouge!8, minimum width=1.5cm, minimum height=8.7cm] at (16.55,4.65) {};
  \node[rotate=90, font=\footnotesize\bfseries, text=insarouge] at (16.55,4.65) {Déploiement et exploitation (chap.~\ref{chap:deploiement})};
  \node[callout, minimum width=16.4cm] at (8.5,-0.55)
    {Des services indépendants, une identité commune, une infrastructure partagée : la plateforme tient à ces trois choix.};
\end{tikzpicture}}
\caption{Les services d'OpenStack étudiés dans le cours. En haut, les services d'exploitation et de plateforme (parties~\ref{part:exploiter}
et~\ref{part:plateformes}) ; au milieu, le socle IaaS (partie~\ref{part:socle}), en rouge plein pour le cœur d'un cloud (calcul, réseau,
images) ; en gris, ce qui les supporte.}
\label{fig:concl-carte}
\end{figure}

\begin{table}[!ht]
\centering
\renewcommand{\arraystretch}{1.25}
\small
\begin{tabular}{@{}>{\raggedright\arraybackslash}p{2.0cm}>{\raggedright\arraybackslash}p{3.9cm}>{\raggedright\arraybackslash}p{8.3cm}r@{}}
\toprule
\textbf{Service} & \textbf{Rôle} & \textbf{Processus et composants à connaître} & \textbf{Chap.} \\
\midrule
Keystone & Identité, authentification, catalogue & Service HTTP (WSGI) ; tokens Fernet ; catalogue des points d'accès & \ref{chap:keystone} \\
Nova & Instances (calcul) & \texttt{nova-api}, \texttt{nova-conductor}, \texttt{nova-scheduler}, \texttt{nova-compute} & \ref{chap:nova} \\
Placement & Inventaire et allocation des ressources & \texttt{placement-api} : fournisseurs, classes, allocations & \ref{chap:placement} \\
Glance & Images & \texttt{glance-api} et ses magasins & \ref{chap:glance} \\
Neutron & Réseau & \texttt{neutron-server}, agents (L3, DHCP), OVS ou OVN & \ref{chap:neutron} \\
Cinder & Volumes (stockage bloc) & \texttt{cinder-api}, \texttt{-scheduler}, \texttt{-volume}, \texttt{-backup} & \ref{chap:cinder} \\
Swift & Stockage objet & Proxy, serveurs de comptes, de conteneurs et d'objets ; anneaux & \ref{chap:swift} \\
Horizon & Tableau de bord & Application Django sous Apache (\texttt{mod\_wsgi}) & \ref{chap:horizon} \\
Octavia & Répartition de charge & \texttt{octavia-api}, worker, health manager, housekeeping ; amphores & \ref{chap:octavia} \\
Heat & Orchestration & \texttt{heat-api}, \texttt{heat-engine} ; modèles HOT & \ref{chap:heat} \\
Ironic & Machines physiques & \texttt{ironic-api}, \texttt{ironic-conductor}, agent \texttt{ironic-python-agent} & \ref{chap:ironic} \\
Magnum & Clusters Kubernetes & \texttt{magnum-api}, \texttt{magnum-conductor} ; pilote de cluster & \ref{chap:magnum} \\
Trove & Bases de données & \texttt{trove-api}, \texttt{trove-taskmanager}, \texttt{trove-conductor} ; agent invité & \ref{chap:trove} \\
Outils & Installer, exploiter, mettre à niveau & Kolla-Ansible, OpenStack-Ansible, DevStack, \dots & \ref{chap:deploiement} \\
\bottomrule
\end{tabular}
\caption{Les services du cours en un coup d'œil.}
\label{tab:concl-services}
\end{table}

% ---------------------------------------------------------------------
\subsection*{Un scénario de bout en bout}
% ---------------------------------------------------------------------

Rien ne montre mieux la coopération des services qu'une demande ordinaire : une machine virtuelle dotée d'un disque de données et
d'une adresse publique. Le chapitre~\ref{chap:presentation} avait suivi la création d'une machine seule (figure~\ref{fig:pres-scenario}) ; ici, cinq
commandes y ajoutent un volume et une adresse publique, et la figure~\ref{fig:concl-seq} en suit les appels.

\par\noindent\begin{minipage}{\linewidth}
\begin{lstlisting}[style=terminal]
$ openstack volume create --size 20 donnees
$ openstack server create --image ubuntu-24.04 --flavor m1.small --network prive vm1
$ openstack server add volume vm1 donnees
$ openstack floating ip create public
$ openstack server add floating ip vm1 203.0.113.10
\end{lstlisting}
\end{minipage}\par

\begin{figure}[!ht]
\centering
\adjustbox{max width=\linewidth}{%
\begin{tikzpicture}[>=Stealth, font=\footnotesize]
  \node[acteur, minimum width=2.4cm] (c) at (0.8,0)  {Client};
  \node[acteur, minimum width=2.4cm, fill=insapeche, draw=insaorange] (k) at (3.9,0)  {Keystone};
  \node[acteur, minimum width=2.4cm, fill=insapeche, draw=insaorange] (ci) at (7.0,0)  {Cinder};
  \node[acteur, minimum width=2.4cm, fill=insarouge, text=white, draw=black] (n) at (10.1,0) {Nova};
  \node[acteur, minimum width=2.4cm, fill=insapeche, draw=insaorange] (g) at (13.2,0) {Glance};
  \node[acteur, minimum width=2.4cm, fill=insapeche, draw=insaorange] (ne) at (16.3,0) {Neutron};
  \foreach \n in {c,k,ci,n,g,ne} \draw[vie] (\n.south) -- ++(0,-10.9);
  \seqmsg{0.8}{3.9}{-1.3}{1}{authentification}
  \seqrep{3.9}{0.8}{-2.0}{}{token}
  \seqmsg{0.8}{7.0}{-2.9}{2}{\texttt{POST /volumes}}
  \seqrep{7.0}{0.8}{-3.6}{}{volume \texttt{available}}
  \seqmsg{0.8}{10.1}{-4.5}{3}{\texttt{POST /servers}}
  \seqrep{10.1}{0.8}{-5.2}{}{\texttt{202 Accepted}}
  \seqmsg{10.1}{13.2}{-6.1}{4}{image}
  \seqmsg{10.1}{16.3}{-7.0}{5}{port réseau}
  \node[callout, font=\scriptsize] at (10.1,-7.8) {instance \texttt{ACTIVE}};
  \seqmsg{0.8}{10.1}{-8.6}{6}{\texttt{add volume}}
  \seqmsg{10.1}{7.0}{-9.5}{7}{attacher le volume}
  \seqmsg{0.8}{16.3}{-10.4}{8}{adresse flottante}
  \node[callout, font=\scriptsize, text width=4.4cm] at (14.6,-2.9)
    {Chaque appel porte le token ; chaque service vérifie les droits et les quotas.};
\end{tikzpicture}}
\caption{Le scénario en huit échanges. Le client s'authentifie (1), crée le volume (2) et la machine (3) ; Nova répond sans attendre,
puis obtient l'image et le port réseau auprès de Glance et de Neutron (4, 5), détaillés dans la figure~\ref{fig:nova-seq}. Le
volume est attaché par Nova auprès de Cinder (6, 7) et l'adresse flottante associée au port par Neutron (8).}
\label{fig:concl-seq}
\end{figure}

Quatre idées se dégagent. Un seul token, émis par Keystone, sert auprès de tous les services. Les services collaborent par leurs API
REST, sans jamais lire la base les uns des autres. Nova répond \texttt{202 Accepted} avant d'avoir terminé : l'état de l'instance
(\texttt{BUILD}, puis \texttt{ACTIVE}) renseigne sur l'avancement. Enfin, quand une machine finit en \texttt{ERROR}, la bonne question est
« quel service a refusé, et pourquoi ? » (quota, image, port, volume) : on lit l'état de la ressource, puis les journaux du service
concerné.

% ---------------------------------------------------------------------
\subsection*{Les principes qui reviennent partout}
% ---------------------------------------------------------------------

Au-delà de leurs différences, les services partagent une même manière de fonctionner (tableau~\ref{tab:concl-principes}). La
reconnaître dans un service inconnu permet de le comprendre vite : on cherche son API, son processus de coordination, ses agents,
sa base, ses états et ses pilotes. Les bibliothèques \gls{oslo} (configuration, messagerie, accès aux bases, politiques) en sont
le ciment commun.

\begin{table}[!ht]
\centering
\renewcommand{\arraystretch}{1.25}
\small
\begin{tabular}{@{}>{\raggedright\arraybackslash}p{3.2cm}>{\raggedright\arraybackslash}p{9.0cm}>{\raggedright\arraybackslash}p{3.1cm}@{}}
\toprule
\textbf{Principe} & \textbf{En pratique} & \textbf{À revoir} \\
\midrule
API REST authentifiée & Chaque appel porte un token de Keystone ; le service le valide et applique sa politique de droits. & Chap.~\ref{chap:keystone} \\
API, coordination, agents & L'API valide et répond ; un processus de coordination (\emph{conductor}, \emph{engine}, \emph{taskmanager}) décide ; des agents exécutent sur les nœuds. & Chap.~\ref{chap:nova}, \ref{chap:cinder}, \ref{chap:heat}, \ref{chap:ironic}, \ref{chap:trove} \\
Messagerie interne & Les processus d'un même service s'appellent par RPC sur la file de messages (RabbitMQ) ; d'un service à l'autre, on passe par l'API REST. & Chap.~\ref{chap:nova}, \ref{chap:deploiement} \\
Une base par service & Chaque service gère son schéma (\texttt{db sync}) ; les autres ne le lisent jamais. & Chap.~\ref{chap:deploiement} \\
Opérations asynchrones & L'API répond vite ; la ressource passe par des états (\texttt{BUILD} $\to$ \texttt{ACTIVE}, \texttt{creating} $\to$ \texttt{available}) et finit en \texttt{ERROR} s'il y a un échec. & Chap.~\ref{chap:nova}, \ref{chap:cinder}, \ref{chap:magnum} \\
Quotas et capacité & Chaque service plafonne les \glspl{quota} par projet ; Nova s'appuie sur Placement pour connaître la capacité réelle. & Chap.~\ref{chap:nova}, \ref{chap:placement} \\
Pilotes et greffons & Le cœur est stable, le variable est dans des pilotes : ML2 (OVS, OVN), backends de Cinder, magasins de Glance, types matériels d'Ironic, pilotes de cluster de Magnum, datastores de Trove. & Chap.~\ref{chap:neutron}, \ref{chap:cinder}, \ref{chap:ironic}, \ref{chap:magnum}, \ref{chap:trove} \\
Versions d'API négociées & Les \emph{microversions} font évoluer l'API sans casser les clients : le client indique la version qu'il attend (Nova, Placement, Cinder, Ironic). & Chap.~\ref{chap:nova}, \ref{chap:placement} \\
\bottomrule
\end{tabular}
\caption{Huit principes communs aux services d'OpenStack.}
\label{tab:concl-principes}
\end{table}

% ---------------------------------------------------------------------
\subsection*{Diagnostiquer : une méthode commune}
% ---------------------------------------------------------------------

Les principes précédents donnent aussi une méthode de diagnostic qui vaut pour tous les services (figure~\ref{fig:concl-diag}). On
part de la ressource en échec, on remonte au service qui la possède, puis on suit la requête dans ses journaux avant d'élargir aux
dépendances et à l'infrastructure.

\begin{figure}[!ht]
\centering
\adjustbox{max width=\linewidth}{%
\begin{tikzpicture}[>=Stealth, font=\footnotesize,
    cmd/.style={font=\scriptsize, text width=2.9cm, align=center}]
  \foreach \x/\t/\c [count=\i] in {1.6/{Lire\\ l'état}/{\texttt{openstack \dots{} show} : \texttt{status}, champ \texttt{fault}},
        4.9/{Trouver le\\ service}/{à qui appartient la ressource ? (\texttt{endpoint list})},
        8.2/{Suivre la\\ requête}/{identifiant \texttt{req-\dots} dans les journaux (\texttt{--debug})},
        11.5/{Vérifier les\\ dépendances}/{quotas, services et agents \texttt{up}, Placement, Glance, Neutron, Cinder},
        14.8/{Vérifier\\ l'infrastructure}/{RabbitMQ, base de données, horloge (NTP), espace disque}}
    { \node[pfoc, minimum width=2.9cm, minimum height=1.1cm] (s\i) at (\x,1.6) {\t};
      \node[num] at (\x-1.25,2.25) {\i};
      \node[cmd] at (\x,0.35) {\c}; }
  \foreach \a/\b in {1/2,2/3,3/4,4/5} \draw[fluxr] (s\a.east) -- (s\b.west);
  \node[callout, minimum width=15.6cm] at (8.2,-1.0)
    {On élargit le champ pas à pas : la cause est presque toujours dans le service propriétaire ou dans l'une de ses dépendances.};
\end{tikzpicture}}
\caption{Cinq temps pour diagnostiquer une ressource en échec, quel que soit le service.}
\label{fig:concl-diag}
\end{figure}

% ---------------------------------------------------------------------
\subsection*{Quel service pour quel besoin}
% ---------------------------------------------------------------------

Le tableau~\ref{tab:concl-besoins} part du besoin, ce qui est la démarche d'un utilisateur ou d'un architecte.

\begin{table}[!ht]
\centering
\renewcommand{\arraystretch}{1.25}
\small
\begin{tabular}{@{}>{\raggedright\arraybackslash}p{4.5cm}>{\raggedright\arraybackslash}p{3.7cm}>{\raggedright\arraybackslash}p{7.2cm}@{}}
\toprule
\textbf{Besoin} & \textbf{Service} & \textbf{À retenir} \\
\midrule
Une machine virtuelle & Nova (avec Glance, Neutron, Placement) & Disque racine local, sauf démarrage sur un volume. \\
Un disque qui survit à la machine & Cinder & Volume attaché à une instance, avec instantanés et sauvegardes. \\
Des fichiers, des archives, via HTTP & Swift & Indépendant des instances ; réplication ou codage d'effacement. \\
Réseau, routeur, adresse flottante & Neutron & Réseaux de projet isolés, groupes de sécurité. \\
Répartir le trafic entre serveurs & Octavia & Répartiteur fourni comme service, rôle des amphores. \\
Recréer une infrastructure à l'identique & Heat & Un modèle HOT décrit une pile entière. \\
Une machine physique & Ironic & Livraison de serveurs nus, souvent derrière Nova. \\
Un cluster Kubernetes & Magnum & Modèle de cluster, pilote Cluster API. \\
Une base de données managée & Trove & Instance, volume et moteur dans un conteneur. \\
Une interface web & Horizon & Appelle les API avec les droits de l'utilisateur. \\
\bottomrule
\end{tabular}
\caption{Du besoin au service.}
\label{tab:concl-besoins}
\end{table}

% ---------------------------------------------------------------------
\subsection*{Mettre en production : une liste de contrôle}
% ---------------------------------------------------------------------

Les chapitres ont chacun leur section d'exploitation ; en voici l'essentiel, dans l'ordre où l'on s'y heurte (chapitre~\ref{chap:deploiement}) :

\begin{itemize}
  \item \textbf{Topologie} : au moins trois contrôleurs, adresse virtuelle devant les API, réseaux de gestion, de tunnel et externe
        séparés.
  \item \textbf{Sécurité} : TLS sur les API exposées, secrets hors des dépôts et renouvelés, réseau de gestion interne.
  \item \textbf{Capacité} : \glspl{quota} par projet, gabarits (\emph{flavors}) et images de référence, surveillance de l'occupation des
        hyperviseurs.
  \item \textbf{Sauvegardes} : base de données de chaque service, configuration de l'outil de déploiement, volumes
        (chapitre~\ref{chap:cinder}) ; une restauration répétée au moins une fois.
  \item \textbf{Supervision} : état des services et des agents, journaux centralisés, alertes.
  \item \textbf{Pannes} : une procédure pour la perte d'un hyperviseur (évacuation) et pour la perte d'un contrôleur.
  \item \textbf{Mises à niveau} : un plan par version (versions \gls{slurp}), un environnement d'essai, une sauvegarde avant de
        commencer.
\end{itemize}

% ---------------------------------------------------------------------
\subsection*{Limites, alternatives et évolutions}
% ---------------------------------------------------------------------

\textbf{Limites.} OpenStack est puissant, mais il se paie en exploitation : de nombreux processus à superviser, un cycle de mise à
niveau tous les six mois, une documentation répartie projet par projet, une courbe d'apprentissage réelle. Cet effort se justifie
quand l'organisation en retire un bénéfice net (voir le tableau~\ref{tab:concl-alternatives}) ; il ne se justifie pas pour quelques
machines.

\begin{table}[!ht]
\centering
\renewcommand{\arraystretch}{1.25}
\small
\begin{tabular}{@{}>{\raggedright\arraybackslash}p{2.4cm}>{\raggedright\arraybackslash}p{4.4cm}>{\raggedright\arraybackslash}p{4.4cm}>{\raggedright\arraybackslash}p{4.1cm}@{}}
\toprule
\textbf{Critère} & \textbf{OpenStack (cloud privé)} & \textbf{Cloud public} & \textbf{Kubernetes seul} \\
\midrule
Ce qui est fourni & Machines, réseaux, volumes, objets, quelques services managés & Idem, avec un catalogue de services managés bien plus large & Conteneurs et leur réseau, sur des machines déjà disponibles \\
Maîtrise & Totale : matériel, données, versions & Encadrée par le contrat du fournisseur & Celle de l'infrastructure sous-jacente \\
Exploitation & À votre charge (chapitre~\ref{chap:deploiement}) & Assurée par le fournisseur & À votre charge, plus légère \\
Élasticité & Bornée par le matériel acheté & Quasi illimitée, à la demande & Bornée par les nœuds \\
Plutôt si\dots & Charge stable, souveraineté, machines physiques, équipe d'exploitation & Charge variable, démarrage rapide, peu d'équipe infrastructure & Applications déjà conteneurisées \\
\bottomrule
\end{tabular}
\caption{Trois façons de se procurer des ressources d'infrastructure (tendances générales, non des règles).}
\label{tab:concl-alternatives}
\end{table}

\textbf{Alternatives.} OpenStack et Kubernetes ne s'excluent pas, ils se complètent : Kubernetes orchestre des conteneurs sur des
machines, OpenStack fournit ces machines, leurs réseaux et leurs disques. Les deux sens de la combinaison existent. Magnum
(chapitre~\ref{chap:magnum}) fournit des clusters Kubernetes \emph{sur} OpenStack ; OpenStack-Helm, Sunbeam et RHOSO déploient OpenStack
\emph{sur} Kubernetes (section~\ref{sec:dep-autres}).

\textbf{Évolutions.} Quelques mouvements sont visibles dans le cours même :
\begin{itemize}
  \item un rythme de publication semestriel, avec des versions \gls{slurp} pour ceux qui ne mettent à niveau qu'une fois par an
        \parencite{openstack-releases,openstack-slurp} ;
  \item le départ de certains outils : TripleO est retiré, Swift n'est plus déployé par Kolla-Ansible (section~\ref{sec:dep-kolla}) ;
  \item le rapprochement avec Kubernetes, par le déploiement (Helm, Sunbeam, RHOSO) et par les services (pilotes Cluster API de Magnum,
        opérateurs de bases de données comme alternative à Trove) ;
  \item un cadre institutionnel nouveau : l'OpenInfra Foundation a rejoint la Linux Foundation le 3~juin 2025, tandis que
        OpenStack garde sa gouvernance communautaire \parencite{openinfra-lf}.
\end{itemize}

% ---------------------------------------------------------------------
\subsection*{Ce que le cours n'a pas couvert}
% ---------------------------------------------------------------------

OpenStack compte bien plus de projets officiels que les quinze chapitres \parencite{governance-projects}. Le tableau~\ref{tab:concl-autres}
en cite quelques-uns, qui complètent naturellement les services vus.

\begin{table}[!ht]
\centering
\renewcommand{\arraystretch}{1.2}
\small
\begin{tabular}{@{}>{\raggedright\arraybackslash}p{2.5cm}>{\raggedright\arraybackslash}p{5.8cm}>{\raggedright\arraybackslash}p{7.1cm}@{}}
\toprule
\textbf{Projet} & \textbf{Rôle (selon le registre officiel)} & \textbf{Lien avec le cours} \\
\midrule
Barbican & Gestionnaire de clés & Certificats des clusters de Magnum \\
Designate & Service DNS & Noms de domaine pour les instances et les adresses flottantes \\
Manila & Systèmes de fichiers partagés & Complète Cinder (blocs) et Swift (objets) \\
Skyline & Tableau de bord & Alternative à Horizon \\
Telemetry & Service de télémétrie & Mesures de l'usage des ressources, base de la supervision et de la facturation \\
Masakari & Haute disponibilité des instances & Relance des instances d'un hôte tombé (complète l'évacuation manuelle) \\
Watcher & Optimisation de l'infrastructure & Répartition de la charge entre hyperviseurs \\
Blazar & Réservation de ressources & Réserver des machines pour une période \\
Cyborg & Gestion du cycle de vie des accélérateurs & Accélérateurs matériels (GPU, FPGA) mis à disposition des instances \\
Tacker & Orchestration NFV & Fonctions réseau virtualisées, notamment chez les opérateurs de télécoms \\
CloudKitty & Service de tarification (\emph{rating}) & Facturation à partir des mesures \\
Zun & Service de conteneurs & Lancer des conteneurs sans passer par un cluster Kubernetes \\
\bottomrule
\end{tabular}
\caption{Quelques projets officiels non étudiés. Les rôles reprennent les descriptions du registre ; la colonne de droite est une lecture
d'enseignement, non une affirmation du registre.}
\label{tab:concl-autres}
\end{table}

% ---------------------------------------------------------------------
\subsection*{Pour aller plus loin}
% ---------------------------------------------------------------------

\begin{itemize}
  \item \textbf{Pratiquer} : installer DevStack ou Sunbeam sur une machine jetable (chapitre~\ref{chap:deploiement}), puis refaire le
        scénario de la figure~\ref{fig:concl-seq}, écrire un modèle Heat, ouvrir un répartiteur Octavia. Lire un journal de service
        vaut tous les schémas.
  \item \textbf{Lire} : la documentation officielle de chaque projet, le guide d'installation et le guide de haute disponibilité
        \parencite{install-guide,ha-guide}.
  \item \textbf{Se certifier} : l'examen COA (\emph{Certified OpenStack Administrator}) dure trois heures et se passe en ligne,
        sur des tâches réalisées en ligne de commande et dans Horizon ; sa page indique qu'il repose sur la version Caracal, plus
        ancienne que celles du cours, et conseille six mois d'expérience préalable \parencite{coa}.
  \item \textbf{Contribuer} : le guide du contributeur couvre le code, la documentation, les traductions et les retours
        d'opérateurs \parencite{openstack-contributors}.
\end{itemize}

% ---------------------------------------------------------------------
\subsection*{Messages clés}
% ---------------------------------------------------------------------

\begin{aretenir}
\begin{itemize}
  \item OpenStack est une famille de services indépendants, qui partagent une identité (Keystone), une infrastructure (base,
        messagerie, cache) et des conventions communes.
  \item Une action de l'utilisateur traverse plusieurs services par leurs API REST ; elle est asynchrone, et ses états disent où
        elle en est.
  \item Chaque service suit le même plan : API, coordination, agents, base, pilotes ; comprendre l'un aide à comprendre les autres.
  \item Le choix de l'outil de déploiement structure l'exploitation : Kolla-Ansible ou OpenStack-Ansible en production, DevStack pour
        apprendre et développer.
  \item Un cloud en production demande de l'exploitation : haute disponibilité, TLS, sauvegardes, supervision, mises à niveau.
  \item OpenStack complète le cloud public et Kubernetes plus qu'il ne les remplace : le bon choix dépend de la charge, de la
        maîtrise voulue et de l'équipe disponible.
\end{itemize}
\end{aretenir}

Reste à pratiquer : un cloud se comprend d'abord en le faisant tourner.
.
%  Section non numérotée, ajoutée au sommaire.
%  Sources : les chapitres du cours, le registre des projets officiels
%  (gouvernance d'OpenStack), l'annonce de l'arrivée de l'OpenInfra Foundation
%  à la Linux Foundation, la page de l'examen COA et le guide du contributeur
%  (voir sources.bib). Conclusion détaillée : 10 pages au plus.
% =====================================================================

\phantomsection
\section*{Conclusion}
\addcontentsline{toc}{section}{Conclusion}

% Figures et tableaux de cette section non numérotée : numérotation propre (C.1, C.2\dots)
% pour ne pas prolonger celle du chapitre précédent.
\setcounter{figure}{0}\setcounter{table}{0}
\renewcommand{\thefigure}{C.\arabic{figure}}\renewcommand{\thetable}{C.\arabic{table}}
\renewcommand{\theHfigure}{C.\arabic{figure}}\renewcommand{\theHtable}{C.\arabic{table}}

Ce cours a présenté OpenStack comme une famille de services indépendants qui, ensemble, transforment des serveurs, des disques et
des réseaux en un cloud en libre-service. Il est parti de la carte d'ensemble et de l'identité (Keystone, partie~\ref{part:fondations}),
a suivi la vie d'une machine virtuelle à travers le socle IaaS (Nova, Placement, Glance, Neutron, Cinder, Swift,
partie~\ref{part:socle}), puis a présenté les services qui permettent d'exploiter le cloud (Horizon, Octavia, Heat,
partie~\ref{part:exploiter}) et ceux qui fournissent d'autres ressources au-dessus du socle (Ironic, Magnum, Trove,
partie~\ref{part:plateformes}). Il s'est achevé du côté de l'administrateur, avec le déploiement (partie~\ref{part:miseenoeuvre}). Cette conclusion rassemble l'ensemble : une carte,
un scénario de bout en bout, les principes qui reviennent d'un service à l'autre, des repères pour choisir et exploiter, puis les
limites de la plateforme, ce que le cours n'a pas traité et la manière de poursuivre.

% ---------------------------------------------------------------------
\subsection*{La carte d'ensemble}
% ---------------------------------------------------------------------

La figure~\ref{fig:concl-carte} reprend la carte du chapitre~\ref{chap:presentation} (figure~\ref{fig:pres-archi}) en y ajoutant les
numéros de chapitre. Tous les services reposent sur la même infrastructure (base de données, file de
messages, cache) et sur Keystone, qui authentifie chaque appel ; l'outil de déploiement installe et met à jour le tout. Le
tableau~\ref{tab:concl-services} rappelle, pour chaque service, son rôle et les processus à connaître.

\begin{figure}[!ht]
\centering
\adjustbox{max width=\linewidth}{%
\begin{tikzpicture}[>=Stealth, font=\footnotesize]
  \node[pcli, minimum width=13.0cm] at (8.5,8.3)
    {Utilisateurs \ $\cdot$ \ ligne de commande \ $\cdot$ \ API REST \ $\cdot$ \ Horizon (chap.~\ref{chap:horizon})};
  % services d'exploitation et de plateforme
  \node[psvc, minimum width=2.6cm, minimum height=1.1cm, inner xsep=2pt] at (2.7,6.8) {\textbf{Heat} {\scriptsize\mdseries(chap.~\ref{chap:heat})}\\ {\scriptsize\mdseries orchestration}};
  \node[psvc, minimum width=2.6cm, minimum height=1.1cm, inner xsep=2pt] at (5.6,6.8) {\textbf{Octavia} {\scriptsize\mdseries(chap.~\ref{chap:octavia})}\\ {\scriptsize\mdseries répartition de charge}};
  \node[psvc, minimum width=2.6cm, minimum height=1.1cm, inner xsep=2pt] at (8.5,6.8) {\textbf{Ironic} {\scriptsize\mdseries(chap.~\ref{chap:ironic})}\\ {\scriptsize\mdseries machines physiques}};
  \node[psvc, minimum width=2.6cm, minimum height=1.1cm, inner xsep=2pt] at (11.4,6.8) {\textbf{Magnum} {\scriptsize\mdseries(chap.~\ref{chap:magnum})}\\ {\scriptsize\mdseries Kubernetes managé}};
  \node[psvc, minimum width=2.6cm, minimum height=1.1cm, inner xsep=2pt] at (14.3,6.8) {\textbf{Trove} {\scriptsize\mdseries(chap.~\ref{chap:trove})}\\ {\scriptsize\mdseries bases de données}};
  % services de base
  \node[pfoc, minimum width=3.75cm, minimum height=1.1cm] at (4.5,5.2) {Nova, Placement\\ {\scriptsize\mdseries calcul (chap.~\ref{chap:nova}, \ref{chap:placement})}};
  \node[pfoc, minimum width=3.75cm, minimum height=1.1cm] at (8.5,5.2) {Neutron\\ {\scriptsize\mdseries réseau (chap.~\ref{chap:neutron})}};
  \node[pfoc, minimum width=3.75cm, minimum height=1.1cm] at (12.5,5.2) {Glance\\ {\scriptsize\mdseries images (chap.~\ref{chap:glance})}};
  \node[psvc, minimum width=3.75cm, minimum height=1.1cm] at (6.5,3.75) {Cinder\\ {\scriptsize\mdseries volumes (chap.~\ref{chap:cinder})}};
  \node[psvc, minimum width=3.75cm, minimum height=1.1cm] at (10.5,3.75) {Swift\\ {\scriptsize\mdseries objets (chap.~\ref{chap:swift})}};
  % infrastructure
  \node[pgris, minimum width=13.0cm, minimum height=0.95cm] at (8.5,2.3)
    {Infrastructure partagée\\ {\scriptsize\mdseries MariaDB \ $\cdot$ \ RabbitMQ \ $\cdot$ \ memcached \ $\cdot$ \ HAProxy}};
  \node[pgris, minimum width=13.0cm, minimum height=0.95cm] at (8.5,1.05)
    {Machines et réseaux\\ {\scriptsize\mdseries hyperviseurs, disques, réseau physique, serveurs nus}};
  % barres latérales
  \node[draw=insarouge, thick, rounded corners=4pt, fill=insarouge!8, minimum width=1.5cm, minimum height=8.7cm] at (0.45,4.65) {};
  \node[rotate=90, font=\footnotesize\bfseries, text=insarouge] at (0.45,4.65) {Keystone : identité, catalogue, tokens (chap.~\ref{chap:keystone})};
  \node[draw=insarouge, thick, rounded corners=4pt, fill=insarouge!8, minimum width=1.5cm, minimum height=8.7cm] at (16.55,4.65) {};
  \node[rotate=90, font=\footnotesize\bfseries, text=insarouge] at (16.55,4.65) {Déploiement et exploitation (chap.~\ref{chap:deploiement})};
  \node[callout, minimum width=16.4cm] at (8.5,-0.55)
    {Des services indépendants, une identité commune, une infrastructure partagée : la plateforme tient à ces trois choix.};
\end{tikzpicture}}
\caption{Les services d'OpenStack étudiés dans le cours. En haut, les services d'exploitation et de plateforme (parties~\ref{part:exploiter}
et~\ref{part:plateformes}) ; au milieu, le socle IaaS (partie~\ref{part:socle}), en rouge plein pour le cœur d'un cloud (calcul, réseau,
images) ; en gris, ce qui les supporte.}
\label{fig:concl-carte}
\end{figure}

\begin{table}[!ht]
\centering
\renewcommand{\arraystretch}{1.25}
\small
\begin{tabular}{@{}>{\raggedright\arraybackslash}p{2.0cm}>{\raggedright\arraybackslash}p{3.9cm}>{\raggedright\arraybackslash}p{8.3cm}r@{}}
\toprule
\textbf{Service} & \textbf{Rôle} & \textbf{Processus et composants à connaître} & \textbf{Chap.} \\
\midrule
Keystone & Identité, authentification, catalogue & Service HTTP (WSGI) ; tokens Fernet ; catalogue des points d'accès & \ref{chap:keystone} \\
Nova & Instances (calcul) & \texttt{nova-api}, \texttt{nova-conductor}, \texttt{nova-scheduler}, \texttt{nova-compute} & \ref{chap:nova} \\
Placement & Inventaire et allocation des ressources & \texttt{placement-api} : fournisseurs, classes, allocations & \ref{chap:placement} \\
Glance & Images & \texttt{glance-api} et ses magasins & \ref{chap:glance} \\
Neutron & Réseau & \texttt{neutron-server}, agents (L3, DHCP), OVS ou OVN & \ref{chap:neutron} \\
Cinder & Volumes (stockage bloc) & \texttt{cinder-api}, \texttt{-scheduler}, \texttt{-volume}, \texttt{-backup} & \ref{chap:cinder} \\
Swift & Stockage objet & Proxy, serveurs de comptes, de conteneurs et d'objets ; anneaux & \ref{chap:swift} \\
Horizon & Tableau de bord & Application Django sous Apache (\texttt{mod\_wsgi}) & \ref{chap:horizon} \\
Octavia & Répartition de charge & \texttt{octavia-api}, worker, health manager, housekeeping ; amphores & \ref{chap:octavia} \\
Heat & Orchestration & \texttt{heat-api}, \texttt{heat-engine} ; modèles HOT & \ref{chap:heat} \\
Ironic & Machines physiques & \texttt{ironic-api}, \texttt{ironic-conductor}, agent \texttt{ironic-python-agent} & \ref{chap:ironic} \\
Magnum & Clusters Kubernetes & \texttt{magnum-api}, \texttt{magnum-conductor} ; pilote de cluster & \ref{chap:magnum} \\
Trove & Bases de données & \texttt{trove-api}, \texttt{trove-taskmanager}, \texttt{trove-conductor} ; agent invité & \ref{chap:trove} \\
Outils & Installer, exploiter, mettre à niveau & Kolla-Ansible, OpenStack-Ansible, DevStack, \dots & \ref{chap:deploiement} \\
\bottomrule
\end{tabular}
\caption{Les services du cours en un coup d'œil.}
\label{tab:concl-services}
\end{table}

% ---------------------------------------------------------------------
\subsection*{Un scénario de bout en bout}
% ---------------------------------------------------------------------

Rien ne montre mieux la coopération des services qu'une demande ordinaire : une machine virtuelle dotée d'un disque de données et
d'une adresse publique. Le chapitre~\ref{chap:presentation} avait suivi la création d'une machine seule (figure~\ref{fig:pres-scenario}) ; ici, cinq
commandes y ajoutent un volume et une adresse publique, et la figure~\ref{fig:concl-seq} en suit les appels.

\par\noindent\begin{minipage}{\linewidth}
\begin{lstlisting}[style=terminal]
$ openstack volume create --size 20 donnees
$ openstack server create --image ubuntu-24.04 --flavor m1.small --network prive vm1
$ openstack server add volume vm1 donnees
$ openstack floating ip create public
$ openstack server add floating ip vm1 203.0.113.10
\end{lstlisting}
\end{minipage}\par

\begin{figure}[!ht]
\centering
\adjustbox{max width=\linewidth}{%
\begin{tikzpicture}[>=Stealth, font=\footnotesize]
  \node[acteur, minimum width=2.4cm] (c) at (0.8,0)  {Client};
  \node[acteur, minimum width=2.4cm, fill=insapeche, draw=insaorange] (k) at (3.9,0)  {Keystone};
  \node[acteur, minimum width=2.4cm, fill=insapeche, draw=insaorange] (ci) at (7.0,0)  {Cinder};
  \node[acteur, minimum width=2.4cm, fill=insarouge, text=white, draw=black] (n) at (10.1,0) {Nova};
  \node[acteur, minimum width=2.4cm, fill=insapeche, draw=insaorange] (g) at (13.2,0) {Glance};
  \node[acteur, minimum width=2.4cm, fill=insapeche, draw=insaorange] (ne) at (16.3,0) {Neutron};
  \foreach \n in {c,k,ci,n,g,ne} \draw[vie] (\n.south) -- ++(0,-10.9);
  \seqmsg{0.8}{3.9}{-1.3}{1}{authentification}
  \seqrep{3.9}{0.8}{-2.0}{}{token}
  \seqmsg{0.8}{7.0}{-2.9}{2}{\texttt{POST /volumes}}
  \seqrep{7.0}{0.8}{-3.6}{}{volume \texttt{available}}
  \seqmsg{0.8}{10.1}{-4.5}{3}{\texttt{POST /servers}}
  \seqrep{10.1}{0.8}{-5.2}{}{\texttt{202 Accepted}}
  \seqmsg{10.1}{13.2}{-6.1}{4}{image}
  \seqmsg{10.1}{16.3}{-7.0}{5}{port réseau}
  \node[callout, font=\scriptsize] at (10.1,-7.8) {instance \texttt{ACTIVE}};
  \seqmsg{0.8}{10.1}{-8.6}{6}{\texttt{add volume}}
  \seqmsg{10.1}{7.0}{-9.5}{7}{attacher le volume}
  \seqmsg{0.8}{16.3}{-10.4}{8}{adresse flottante}
  \node[callout, font=\scriptsize, text width=4.4cm] at (14.6,-2.9)
    {Chaque appel porte le token ; chaque service vérifie les droits et les quotas.};
\end{tikzpicture}}
\caption{Le scénario en huit échanges. Le client s'authentifie (1), crée le volume (2) et la machine (3) ; Nova répond sans attendre,
puis obtient l'image et le port réseau auprès de Glance et de Neutron (4, 5), détaillés dans la figure~\ref{fig:nova-seq}. Le
volume est attaché par Nova auprès de Cinder (6, 7) et l'adresse flottante associée au port par Neutron (8).}
\label{fig:concl-seq}
\end{figure}

Quatre idées se dégagent. Un seul token, émis par Keystone, sert auprès de tous les services. Les services collaborent par leurs API
REST, sans jamais lire la base les uns des autres. Nova répond \texttt{202 Accepted} avant d'avoir terminé : l'état de l'instance
(\texttt{BUILD}, puis \texttt{ACTIVE}) renseigne sur l'avancement. Enfin, quand une machine finit en \texttt{ERROR}, la bonne question est
« quel service a refusé, et pourquoi ? » (quota, image, port, volume) : on lit l'état de la ressource, puis les journaux du service
concerné.

% ---------------------------------------------------------------------
\subsection*{Les principes qui reviennent partout}
% ---------------------------------------------------------------------

Au-delà de leurs différences, les services partagent une même manière de fonctionner (tableau~\ref{tab:concl-principes}). La
reconnaître dans un service inconnu permet de le comprendre vite : on cherche son API, son processus de coordination, ses agents,
sa base, ses états et ses pilotes. Les bibliothèques \gls{oslo} (configuration, messagerie, accès aux bases, politiques) en sont
le ciment commun.

\begin{table}[!ht]
\centering
\renewcommand{\arraystretch}{1.25}
\small
\begin{tabular}{@{}>{\raggedright\arraybackslash}p{3.2cm}>{\raggedright\arraybackslash}p{9.0cm}>{\raggedright\arraybackslash}p{3.1cm}@{}}
\toprule
\textbf{Principe} & \textbf{En pratique} & \textbf{À revoir} \\
\midrule
API REST authentifiée & Chaque appel porte un token de Keystone ; le service le valide et applique sa politique de droits. & Chap.~\ref{chap:keystone} \\
API, coordination, agents & L'API valide et répond ; un processus de coordination (\emph{conductor}, \emph{engine}, \emph{taskmanager}) décide ; des agents exécutent sur les nœuds. & Chap.~\ref{chap:nova}, \ref{chap:cinder}, \ref{chap:heat}, \ref{chap:ironic}, \ref{chap:trove} \\
Messagerie interne & Les processus d'un même service s'appellent par RPC sur la file de messages (RabbitMQ) ; d'un service à l'autre, on passe par l'API REST. & Chap.~\ref{chap:nova}, \ref{chap:deploiement} \\
Une base par service & Chaque service gère son schéma (\texttt{db sync}) ; les autres ne le lisent jamais. & Chap.~\ref{chap:deploiement} \\
Opérations asynchrones & L'API répond vite ; la ressource passe par des états (\texttt{BUILD} $\to$ \texttt{ACTIVE}, \texttt{creating} $\to$ \texttt{available}) et finit en \texttt{ERROR} s'il y a un échec. & Chap.~\ref{chap:nova}, \ref{chap:cinder}, \ref{chap:magnum} \\
Quotas et capacité & Chaque service plafonne les \glspl{quota} par projet ; Nova s'appuie sur Placement pour connaître la capacité réelle. & Chap.~\ref{chap:nova}, \ref{chap:placement} \\
Pilotes et greffons & Le cœur est stable, le variable est dans des pilotes : ML2 (OVS, OVN), backends de Cinder, magasins de Glance, types matériels d'Ironic, pilotes de cluster de Magnum, datastores de Trove. & Chap.~\ref{chap:neutron}, \ref{chap:cinder}, \ref{chap:ironic}, \ref{chap:magnum}, \ref{chap:trove} \\
Versions d'API négociées & Les \emph{microversions} font évoluer l'API sans casser les clients : le client indique la version qu'il attend (Nova, Placement, Cinder, Ironic). & Chap.~\ref{chap:nova}, \ref{chap:placement} \\
\bottomrule
\end{tabular}
\caption{Huit principes communs aux services d'OpenStack.}
\label{tab:concl-principes}
\end{table}

% ---------------------------------------------------------------------
\subsection*{Diagnostiquer : une méthode commune}
% ---------------------------------------------------------------------

Les principes précédents donnent aussi une méthode de diagnostic qui vaut pour tous les services (figure~\ref{fig:concl-diag}). On
part de la ressource en échec, on remonte au service qui la possède, puis on suit la requête dans ses journaux avant d'élargir aux
dépendances et à l'infrastructure.

\begin{figure}[!ht]
\centering
\adjustbox{max width=\linewidth}{%
\begin{tikzpicture}[>=Stealth, font=\footnotesize,
    cmd/.style={font=\scriptsize, text width=2.9cm, align=center}]
  \foreach \x/\t/\c [count=\i] in {1.6/{Lire\\ l'état}/{\texttt{openstack \dots{} show} : \texttt{status}, champ \texttt{fault}},
        4.9/{Trouver le\\ service}/{à qui appartient la ressource ? (\texttt{endpoint list})},
        8.2/{Suivre la\\ requête}/{identifiant \texttt{req-\dots} dans les journaux (\texttt{--debug})},
        11.5/{Vérifier les\\ dépendances}/{quotas, services et agents \texttt{up}, Placement, Glance, Neutron, Cinder},
        14.8/{Vérifier\\ l'infrastructure}/{RabbitMQ, base de données, horloge (NTP), espace disque}}
    { \node[pfoc, minimum width=2.9cm, minimum height=1.1cm] (s\i) at (\x,1.6) {\t};
      \node[num] at (\x-1.25,2.25) {\i};
      \node[cmd] at (\x,0.35) {\c}; }
  \foreach \a/\b in {1/2,2/3,3/4,4/5} \draw[fluxr] (s\a.east) -- (s\b.west);
  \node[callout, minimum width=15.6cm] at (8.2,-1.0)
    {On élargit le champ pas à pas : la cause est presque toujours dans le service propriétaire ou dans l'une de ses dépendances.};
\end{tikzpicture}}
\caption{Cinq temps pour diagnostiquer une ressource en échec, quel que soit le service.}
\label{fig:concl-diag}
\end{figure}

% ---------------------------------------------------------------------
\subsection*{Quel service pour quel besoin}
% ---------------------------------------------------------------------

Le tableau~\ref{tab:concl-besoins} part du besoin, ce qui est la démarche d'un utilisateur ou d'un architecte.

\begin{table}[!ht]
\centering
\renewcommand{\arraystretch}{1.25}
\small
\begin{tabular}{@{}>{\raggedright\arraybackslash}p{4.5cm}>{\raggedright\arraybackslash}p{3.7cm}>{\raggedright\arraybackslash}p{7.2cm}@{}}
\toprule
\textbf{Besoin} & \textbf{Service} & \textbf{À retenir} \\
\midrule
Une machine virtuelle & Nova (avec Glance, Neutron, Placement) & Disque racine local, sauf démarrage sur un volume. \\
Un disque qui survit à la machine & Cinder & Volume attaché à une instance, avec instantanés et sauvegardes. \\
Des fichiers, des archives, via HTTP & Swift & Indépendant des instances ; réplication ou codage d'effacement. \\
Réseau, routeur, adresse flottante & Neutron & Réseaux de projet isolés, groupes de sécurité. \\
Répartir le trafic entre serveurs & Octavia & Répartiteur fourni comme service, rôle des amphores. \\
Recréer une infrastructure à l'identique & Heat & Un modèle HOT décrit une pile entière. \\
Une machine physique & Ironic & Livraison de serveurs nus, souvent derrière Nova. \\
Un cluster Kubernetes & Magnum & Modèle de cluster, pilote Cluster API. \\
Une base de données managée & Trove & Instance, volume et moteur dans un conteneur. \\
Une interface web & Horizon & Appelle les API avec les droits de l'utilisateur. \\
\bottomrule
\end{tabular}
\caption{Du besoin au service.}
\label{tab:concl-besoins}
\end{table}

% ---------------------------------------------------------------------
\subsection*{Mettre en production : une liste de contrôle}
% ---------------------------------------------------------------------

Les chapitres ont chacun leur section d'exploitation ; en voici l'essentiel, dans l'ordre où l'on s'y heurte (chapitre~\ref{chap:deploiement}) :

\begin{itemize}
  \item \textbf{Topologie} : au moins trois contrôleurs, adresse virtuelle devant les API, réseaux de gestion, de tunnel et externe
        séparés.
  \item \textbf{Sécurité} : TLS sur les API exposées, secrets hors des dépôts et renouvelés, réseau de gestion interne.
  \item \textbf{Capacité} : \glspl{quota} par projet, gabarits (\emph{flavors}) et images de référence, surveillance de l'occupation des
        hyperviseurs.
  \item \textbf{Sauvegardes} : base de données de chaque service, configuration de l'outil de déploiement, volumes
        (chapitre~\ref{chap:cinder}) ; une restauration répétée au moins une fois.
  \item \textbf{Supervision} : état des services et des agents, journaux centralisés, alertes.
  \item \textbf{Pannes} : une procédure pour la perte d'un hyperviseur (évacuation) et pour la perte d'un contrôleur.
  \item \textbf{Mises à niveau} : un plan par version (versions \gls{slurp}), un environnement d'essai, une sauvegarde avant de
        commencer.
\end{itemize}

% ---------------------------------------------------------------------
\subsection*{Limites, alternatives et évolutions}
% ---------------------------------------------------------------------

\textbf{Limites.} OpenStack est puissant, mais il se paie en exploitation : de nombreux processus à superviser, un cycle de mise à
niveau tous les six mois, une documentation répartie projet par projet, une courbe d'apprentissage réelle. Cet effort se justifie
quand l'organisation en retire un bénéfice net (voir le tableau~\ref{tab:concl-alternatives}) ; il ne se justifie pas pour quelques
machines.

\begin{table}[!ht]
\centering
\renewcommand{\arraystretch}{1.25}
\small
\begin{tabular}{@{}>{\raggedright\arraybackslash}p{2.4cm}>{\raggedright\arraybackslash}p{4.4cm}>{\raggedright\arraybackslash}p{4.4cm}>{\raggedright\arraybackslash}p{4.1cm}@{}}
\toprule
\textbf{Critère} & \textbf{OpenStack (cloud privé)} & \textbf{Cloud public} & \textbf{Kubernetes seul} \\
\midrule
Ce qui est fourni & Machines, réseaux, volumes, objets, quelques services managés & Idem, avec un catalogue de services managés bien plus large & Conteneurs et leur réseau, sur des machines déjà disponibles \\
Maîtrise & Totale : matériel, données, versions & Encadrée par le contrat du fournisseur & Celle de l'infrastructure sous-jacente \\
Exploitation & À votre charge (chapitre~\ref{chap:deploiement}) & Assurée par le fournisseur & À votre charge, plus légère \\
Élasticité & Bornée par le matériel acheté & Quasi illimitée, à la demande & Bornée par les nœuds \\
Plutôt si\dots & Charge stable, souveraineté, machines physiques, équipe d'exploitation & Charge variable, démarrage rapide, peu d'équipe infrastructure & Applications déjà conteneurisées \\
\bottomrule
\end{tabular}
\caption{Trois façons de se procurer des ressources d'infrastructure (tendances générales, non des règles).}
\label{tab:concl-alternatives}
\end{table}

\textbf{Alternatives.} OpenStack et Kubernetes ne s'excluent pas, ils se complètent : Kubernetes orchestre des conteneurs sur des
machines, OpenStack fournit ces machines, leurs réseaux et leurs disques. Les deux sens de la combinaison existent. Magnum
(chapitre~\ref{chap:magnum}) fournit des clusters Kubernetes \emph{sur} OpenStack ; OpenStack-Helm, Sunbeam et RHOSO déploient OpenStack
\emph{sur} Kubernetes (section~\ref{sec:dep-autres}).

\textbf{Évolutions.} Quelques mouvements sont visibles dans le cours même :
\begin{itemize}
  \item un rythme de publication semestriel, avec des versions \gls{slurp} pour ceux qui ne mettent à niveau qu'une fois par an
        \parencite{openstack-releases,openstack-slurp} ;
  \item le départ de certains outils : TripleO est retiré, Swift n'est plus déployé par Kolla-Ansible (section~\ref{sec:dep-kolla}) ;
  \item le rapprochement avec Kubernetes, par le déploiement (Helm, Sunbeam, RHOSO) et par les services (pilotes Cluster API de Magnum,
        opérateurs de bases de données comme alternative à Trove) ;
  \item un cadre institutionnel nouveau : l'OpenInfra Foundation a rejoint la Linux Foundation le 3~juin 2025, tandis que
        OpenStack garde sa gouvernance communautaire \parencite{openinfra-lf}.
\end{itemize}

% ---------------------------------------------------------------------
\subsection*{Ce que le cours n'a pas couvert}
% ---------------------------------------------------------------------

OpenStack compte bien plus de projets officiels que les quinze chapitres \parencite{governance-projects}. Le tableau~\ref{tab:concl-autres}
en cite quelques-uns, qui complètent naturellement les services vus.

\begin{table}[!ht]
\centering
\renewcommand{\arraystretch}{1.2}
\small
\begin{tabular}{@{}>{\raggedright\arraybackslash}p{2.5cm}>{\raggedright\arraybackslash}p{5.8cm}>{\raggedright\arraybackslash}p{7.1cm}@{}}
\toprule
\textbf{Projet} & \textbf{Rôle (selon le registre officiel)} & \textbf{Lien avec le cours} \\
\midrule
Barbican & Gestionnaire de clés & Certificats des clusters de Magnum \\
Designate & Service DNS & Noms de domaine pour les instances et les adresses flottantes \\
Manila & Systèmes de fichiers partagés & Complète Cinder (blocs) et Swift (objets) \\
Skyline & Tableau de bord & Alternative à Horizon \\
Telemetry & Service de télémétrie & Mesures de l'usage des ressources, base de la supervision et de la facturation \\
Masakari & Haute disponibilité des instances & Relance des instances d'un hôte tombé (complète l'évacuation manuelle) \\
Watcher & Optimisation de l'infrastructure & Répartition de la charge entre hyperviseurs \\
Blazar & Réservation de ressources & Réserver des machines pour une période \\
Cyborg & Gestion du cycle de vie des accélérateurs & Accélérateurs matériels (GPU, FPGA) mis à disposition des instances \\
Tacker & Orchestration NFV & Fonctions réseau virtualisées, notamment chez les opérateurs de télécoms \\
CloudKitty & Service de tarification (\emph{rating}) & Facturation à partir des mesures \\
Zun & Service de conteneurs & Lancer des conteneurs sans passer par un cluster Kubernetes \\
\bottomrule
\end{tabular}
\caption{Quelques projets officiels non étudiés. Les rôles reprennent les descriptions du registre ; la colonne de droite est une lecture
d'enseignement, non une affirmation du registre.}
\label{tab:concl-autres}
\end{table}

% ---------------------------------------------------------------------
\subsection*{Pour aller plus loin}
% ---------------------------------------------------------------------

\begin{itemize}
  \item \textbf{Pratiquer} : installer DevStack ou Sunbeam sur une machine jetable (chapitre~\ref{chap:deploiement}), puis refaire le
        scénario de la figure~\ref{fig:concl-seq}, écrire un modèle Heat, ouvrir un répartiteur Octavia. Lire un journal de service
        vaut tous les schémas.
  \item \textbf{Lire} : la documentation officielle de chaque projet, le guide d'installation et le guide de haute disponibilité
        \parencite{install-guide,ha-guide}.
  \item \textbf{Se certifier} : l'examen COA (\emph{Certified OpenStack Administrator}) dure trois heures et se passe en ligne,
        sur des tâches réalisées en ligne de commande et dans Horizon ; sa page indique qu'il repose sur la version Caracal, plus
        ancienne que celles du cours, et conseille six mois d'expérience préalable \parencite{coa}.
  \item \textbf{Contribuer} : le guide du contributeur couvre le code, la documentation, les traductions et les retours
        d'opérateurs \parencite{openstack-contributors}.
\end{itemize}

% ---------------------------------------------------------------------
\subsection*{Messages clés}
% ---------------------------------------------------------------------

\begin{aretenir}
\begin{itemize}
  \item OpenStack est une famille de services indépendants, qui partagent une identité (Keystone), une infrastructure (base,
        messagerie, cache) et des conventions communes.
  \item Une action de l'utilisateur traverse plusieurs services par leurs API REST ; elle est asynchrone, et ses états disent où
        elle en est.
  \item Chaque service suit le même plan : API, coordination, agents, base, pilotes ; comprendre l'un aide à comprendre les autres.
  \item Le choix de l'outil de déploiement structure l'exploitation : Kolla-Ansible ou OpenStack-Ansible en production, DevStack pour
        apprendre et développer.
  \item Un cloud en production demande de l'exploitation : haute disponibilité, TLS, sauvegardes, supervision, mises à niveau.
  \item OpenStack complète le cloud public et Kubernetes plus qu'il ne les remplace : le bon choix dépend de la charge, de la
        maîtrise voulue et de l'équipe disponible.
\end{itemize}
\end{aretenir}

Reste à pratiquer : un cloud se comprend d'abord en le faisant tourner.
.
%  Section non numérotée, ajoutée au sommaire.
%  Sources : les chapitres du cours, le registre des projets officiels
%  (gouvernance d'OpenStack), l'annonce de l'arrivée de l'OpenInfra Foundation
%  à la Linux Foundation, la page de l'examen COA et le guide du contributeur
%  (voir sources.bib). Conclusion détaillée : 10 pages au plus.
% =====================================================================

\phantomsection
\section*{Conclusion}
\addcontentsline{toc}{section}{Conclusion}

% Figures et tableaux de cette section non numérotée : numérotation propre (C.1, C.2\dots)
% pour ne pas prolonger celle du chapitre précédent.
\setcounter{figure}{0}\setcounter{table}{0}
\renewcommand{\thefigure}{C.\arabic{figure}}\renewcommand{\thetable}{C.\arabic{table}}
\renewcommand{\theHfigure}{C.\arabic{figure}}\renewcommand{\theHtable}{C.\arabic{table}}

Ce cours a présenté OpenStack comme une famille de services indépendants qui, ensemble, transforment des serveurs, des disques et
des réseaux en un cloud en libre-service. Il est parti de la carte d'ensemble et de l'identité (Keystone, partie~\ref{part:fondations}),
a suivi la vie d'une machine virtuelle à travers le socle IaaS (Nova, Placement, Glance, Neutron, Cinder, Swift,
partie~\ref{part:socle}), puis a présenté les services qui permettent d'exploiter le cloud (Horizon, Octavia, Heat,
partie~\ref{part:exploiter}) et ceux qui fournissent d'autres ressources au-dessus du socle (Ironic, Magnum, Trove,
partie~\ref{part:plateformes}). Il s'est achevé du côté de l'administrateur, avec le déploiement (partie~\ref{part:miseenoeuvre}). Cette conclusion rassemble l'ensemble : une carte,
un scénario de bout en bout, les principes qui reviennent d'un service à l'autre, des repères pour choisir et exploiter, puis les
limites de la plateforme, ce que le cours n'a pas traité et la manière de poursuivre.

% ---------------------------------------------------------------------
\subsection*{La carte d'ensemble}
% ---------------------------------------------------------------------

La figure~\ref{fig:concl-carte} reprend la carte du chapitre~\ref{chap:presentation} (figure~\ref{fig:pres-archi}) en y ajoutant les
numéros de chapitre. Tous les services reposent sur la même infrastructure (base de données, file de
messages, cache) et sur Keystone, qui authentifie chaque appel ; l'outil de déploiement installe et met à jour le tout. Le
tableau~\ref{tab:concl-services} rappelle, pour chaque service, son rôle et les processus à connaître.

\begin{figure}[!ht]
\centering
\adjustbox{max width=\linewidth}{%
\begin{tikzpicture}[>=Stealth, font=\footnotesize]
  \node[pcli, minimum width=13.0cm] at (8.5,8.3)
    {Utilisateurs \ $\cdot$ \ ligne de commande \ $\cdot$ \ API REST \ $\cdot$ \ Horizon (chap.~\ref{chap:horizon})};
  % services d'exploitation et de plateforme
  \node[psvc, minimum width=2.6cm, minimum height=1.1cm, inner xsep=2pt] at (2.7,6.8) {\textbf{Heat} {\scriptsize\mdseries(chap.~\ref{chap:heat})}\\ {\scriptsize\mdseries orchestration}};
  \node[psvc, minimum width=2.6cm, minimum height=1.1cm, inner xsep=2pt] at (5.6,6.8) {\textbf{Octavia} {\scriptsize\mdseries(chap.~\ref{chap:octavia})}\\ {\scriptsize\mdseries répartition de charge}};
  \node[psvc, minimum width=2.6cm, minimum height=1.1cm, inner xsep=2pt] at (8.5,6.8) {\textbf{Ironic} {\scriptsize\mdseries(chap.~\ref{chap:ironic})}\\ {\scriptsize\mdseries machines physiques}};
  \node[psvc, minimum width=2.6cm, minimum height=1.1cm, inner xsep=2pt] at (11.4,6.8) {\textbf{Magnum} {\scriptsize\mdseries(chap.~\ref{chap:magnum})}\\ {\scriptsize\mdseries Kubernetes managé}};
  \node[psvc, minimum width=2.6cm, minimum height=1.1cm, inner xsep=2pt] at (14.3,6.8) {\textbf{Trove} {\scriptsize\mdseries(chap.~\ref{chap:trove})}\\ {\scriptsize\mdseries bases de données}};
  % services de base
  \node[pfoc, minimum width=3.75cm, minimum height=1.1cm] at (4.5,5.2) {Nova, Placement\\ {\scriptsize\mdseries calcul (chap.~\ref{chap:nova}, \ref{chap:placement})}};
  \node[pfoc, minimum width=3.75cm, minimum height=1.1cm] at (8.5,5.2) {Neutron\\ {\scriptsize\mdseries réseau (chap.~\ref{chap:neutron})}};
  \node[pfoc, minimum width=3.75cm, minimum height=1.1cm] at (12.5,5.2) {Glance\\ {\scriptsize\mdseries images (chap.~\ref{chap:glance})}};
  \node[psvc, minimum width=3.75cm, minimum height=1.1cm] at (6.5,3.75) {Cinder\\ {\scriptsize\mdseries volumes (chap.~\ref{chap:cinder})}};
  \node[psvc, minimum width=3.75cm, minimum height=1.1cm] at (10.5,3.75) {Swift\\ {\scriptsize\mdseries objets (chap.~\ref{chap:swift})}};
  % infrastructure
  \node[pgris, minimum width=13.0cm, minimum height=0.95cm] at (8.5,2.3)
    {Infrastructure partagée\\ {\scriptsize\mdseries MariaDB \ $\cdot$ \ RabbitMQ \ $\cdot$ \ memcached \ $\cdot$ \ HAProxy}};
  \node[pgris, minimum width=13.0cm, minimum height=0.95cm] at (8.5,1.05)
    {Machines et réseaux\\ {\scriptsize\mdseries hyperviseurs, disques, réseau physique, serveurs nus}};
  % barres latérales
  \node[draw=insarouge, thick, rounded corners=4pt, fill=insarouge!8, minimum width=1.5cm, minimum height=8.7cm] at (0.45,4.65) {};
  \node[rotate=90, font=\footnotesize\bfseries, text=insarouge] at (0.45,4.65) {Keystone : identité, catalogue, tokens (chap.~\ref{chap:keystone})};
  \node[draw=insarouge, thick, rounded corners=4pt, fill=insarouge!8, minimum width=1.5cm, minimum height=8.7cm] at (16.55,4.65) {};
  \node[rotate=90, font=\footnotesize\bfseries, text=insarouge] at (16.55,4.65) {Déploiement et exploitation (chap.~\ref{chap:deploiement})};
  \node[callout, minimum width=16.4cm] at (8.5,-0.55)
    {Des services indépendants, une identité commune, une infrastructure partagée : la plateforme tient à ces trois choix.};
\end{tikzpicture}}
\caption{Les services d'OpenStack étudiés dans le cours. En haut, les services d'exploitation et de plateforme (parties~\ref{part:exploiter}
et~\ref{part:plateformes}) ; au milieu, le socle IaaS (partie~\ref{part:socle}), en rouge plein pour le cœur d'un cloud (calcul, réseau,
images) ; en gris, ce qui les supporte.}
\label{fig:concl-carte}
\end{figure}

\begin{table}[!ht]
\centering
\renewcommand{\arraystretch}{1.25}
\small
\begin{tabular}{@{}>{\raggedright\arraybackslash}p{2.0cm}>{\raggedright\arraybackslash}p{3.9cm}>{\raggedright\arraybackslash}p{8.3cm}r@{}}
\toprule
\textbf{Service} & \textbf{Rôle} & \textbf{Processus et composants à connaître} & \textbf{Chap.} \\
\midrule
Keystone & Identité, authentification, catalogue & Service HTTP (WSGI) ; tokens Fernet ; catalogue des points d'accès & \ref{chap:keystone} \\
Nova & Instances (calcul) & \texttt{nova-api}, \texttt{nova-conductor}, \texttt{nova-scheduler}, \texttt{nova-compute} & \ref{chap:nova} \\
Placement & Inventaire et allocation des ressources & \texttt{placement-api} : fournisseurs, classes, allocations & \ref{chap:placement} \\
Glance & Images & \texttt{glance-api} et ses magasins & \ref{chap:glance} \\
Neutron & Réseau & \texttt{neutron-server}, agents (L3, DHCP), OVS ou OVN & \ref{chap:neutron} \\
Cinder & Volumes (stockage bloc) & \texttt{cinder-api}, \texttt{-scheduler}, \texttt{-volume}, \texttt{-backup} & \ref{chap:cinder} \\
Swift & Stockage objet & Proxy, serveurs de comptes, de conteneurs et d'objets ; anneaux & \ref{chap:swift} \\
Horizon & Tableau de bord & Application Django sous Apache (\texttt{mod\_wsgi}) & \ref{chap:horizon} \\
Octavia & Répartition de charge & \texttt{octavia-api}, worker, health manager, housekeeping ; amphores & \ref{chap:octavia} \\
Heat & Orchestration & \texttt{heat-api}, \texttt{heat-engine} ; modèles HOT & \ref{chap:heat} \\
Ironic & Machines physiques & \texttt{ironic-api}, \texttt{ironic-conductor}, agent \texttt{ironic-python-agent} & \ref{chap:ironic} \\
Magnum & Clusters Kubernetes & \texttt{magnum-api}, \texttt{magnum-conductor} ; pilote de cluster & \ref{chap:magnum} \\
Trove & Bases de données & \texttt{trove-api}, \texttt{trove-taskmanager}, \texttt{trove-conductor} ; agent invité & \ref{chap:trove} \\
Outils & Installer, exploiter, mettre à niveau & Kolla-Ansible, OpenStack-Ansible, DevStack, \dots & \ref{chap:deploiement} \\
\bottomrule
\end{tabular}
\caption{Les services du cours en un coup d'œil.}
\label{tab:concl-services}
\end{table}

% ---------------------------------------------------------------------
\subsection*{Un scénario de bout en bout}
% ---------------------------------------------------------------------

Rien ne montre mieux la coopération des services qu'une demande ordinaire : une machine virtuelle dotée d'un disque de données et
d'une adresse publique. Le chapitre~\ref{chap:presentation} avait suivi la création d'une machine seule (figure~\ref{fig:pres-scenario}) ; ici, cinq
commandes y ajoutent un volume et une adresse publique, et la figure~\ref{fig:concl-seq} en suit les appels.

\par\noindent\begin{minipage}{\linewidth}
\begin{lstlisting}[style=terminal]
$ openstack volume create --size 20 donnees
$ openstack server create --image ubuntu-24.04 --flavor m1.small --network prive vm1
$ openstack server add volume vm1 donnees
$ openstack floating ip create public
$ openstack server add floating ip vm1 203.0.113.10
\end{lstlisting}
\end{minipage}\par

\begin{figure}[!ht]
\centering
\adjustbox{max width=\linewidth}{%
\begin{tikzpicture}[>=Stealth, font=\footnotesize]
  \node[acteur, minimum width=2.4cm] (c) at (0.8,0)  {Client};
  \node[acteur, minimum width=2.4cm, fill=insapeche, draw=insaorange] (k) at (3.9,0)  {Keystone};
  \node[acteur, minimum width=2.4cm, fill=insapeche, draw=insaorange] (ci) at (7.0,0)  {Cinder};
  \node[acteur, minimum width=2.4cm, fill=insarouge, text=white, draw=black] (n) at (10.1,0) {Nova};
  \node[acteur, minimum width=2.4cm, fill=insapeche, draw=insaorange] (g) at (13.2,0) {Glance};
  \node[acteur, minimum width=2.4cm, fill=insapeche, draw=insaorange] (ne) at (16.3,0) {Neutron};
  \foreach \n in {c,k,ci,n,g,ne} \draw[vie] (\n.south) -- ++(0,-10.9);
  \seqmsg{0.8}{3.9}{-1.3}{1}{authentification}
  \seqrep{3.9}{0.8}{-2.0}{}{token}
  \seqmsg{0.8}{7.0}{-2.9}{2}{\texttt{POST /volumes}}
  \seqrep{7.0}{0.8}{-3.6}{}{volume \texttt{available}}
  \seqmsg{0.8}{10.1}{-4.5}{3}{\texttt{POST /servers}}
  \seqrep{10.1}{0.8}{-5.2}{}{\texttt{202 Accepted}}
  \seqmsg{10.1}{13.2}{-6.1}{4}{image}
  \seqmsg{10.1}{16.3}{-7.0}{5}{port réseau}
  \node[callout, font=\scriptsize] at (10.1,-7.8) {instance \texttt{ACTIVE}};
  \seqmsg{0.8}{10.1}{-8.6}{6}{\texttt{add volume}}
  \seqmsg{10.1}{7.0}{-9.5}{7}{attacher le volume}
  \seqmsg{0.8}{16.3}{-10.4}{8}{adresse flottante}
  \node[callout, font=\scriptsize, text width=4.4cm] at (14.6,-2.9)
    {Chaque appel porte le token ; chaque service vérifie les droits et les quotas.};
\end{tikzpicture}}
\caption{Le scénario en huit échanges. Le client s'authentifie (1), crée le volume (2) et la machine (3) ; Nova répond sans attendre,
puis obtient l'image et le port réseau auprès de Glance et de Neutron (4, 5), détaillés dans la figure~\ref{fig:nova-seq}. Le
volume est attaché par Nova auprès de Cinder (6, 7) et l'adresse flottante associée au port par Neutron (8).}
\label{fig:concl-seq}
\end{figure}

Quatre idées se dégagent. Un seul token, émis par Keystone, sert auprès de tous les services. Les services collaborent par leurs API
REST, sans jamais lire la base les uns des autres. Nova répond \texttt{202 Accepted} avant d'avoir terminé : l'état de l'instance
(\texttt{BUILD}, puis \texttt{ACTIVE}) renseigne sur l'avancement. Enfin, quand une machine finit en \texttt{ERROR}, la bonne question est
« quel service a refusé, et pourquoi ? » (quota, image, port, volume) : on lit l'état de la ressource, puis les journaux du service
concerné.

% ---------------------------------------------------------------------
\subsection*{Les principes qui reviennent partout}
% ---------------------------------------------------------------------

Au-delà de leurs différences, les services partagent une même manière de fonctionner (tableau~\ref{tab:concl-principes}). La
reconnaître dans un service inconnu permet de le comprendre vite : on cherche son API, son processus de coordination, ses agents,
sa base, ses états et ses pilotes. Les bibliothèques \gls{oslo} (configuration, messagerie, accès aux bases, politiques) en sont
le ciment commun.

\begin{table}[!ht]
\centering
\renewcommand{\arraystretch}{1.25}
\small
\begin{tabular}{@{}>{\raggedright\arraybackslash}p{3.2cm}>{\raggedright\arraybackslash}p{9.0cm}>{\raggedright\arraybackslash}p{3.1cm}@{}}
\toprule
\textbf{Principe} & \textbf{En pratique} & \textbf{À revoir} \\
\midrule
API REST authentifiée & Chaque appel porte un token de Keystone ; le service le valide et applique sa politique de droits. & Chap.~\ref{chap:keystone} \\
API, coordination, agents & L'API valide et répond ; un processus de coordination (\emph{conductor}, \emph{engine}, \emph{taskmanager}) décide ; des agents exécutent sur les nœuds. & Chap.~\ref{chap:nova}, \ref{chap:cinder}, \ref{chap:heat}, \ref{chap:ironic}, \ref{chap:trove} \\
Messagerie interne & Les processus d'un même service s'appellent par RPC sur la file de messages (RabbitMQ) ; d'un service à l'autre, on passe par l'API REST. & Chap.~\ref{chap:nova}, \ref{chap:deploiement} \\
Une base par service & Chaque service gère son schéma (\texttt{db sync}) ; les autres ne le lisent jamais. & Chap.~\ref{chap:deploiement} \\
Opérations asynchrones & L'API répond vite ; la ressource passe par des états (\texttt{BUILD} $\to$ \texttt{ACTIVE}, \texttt{creating} $\to$ \texttt{available}) et finit en \texttt{ERROR} s'il y a un échec. & Chap.~\ref{chap:nova}, \ref{chap:cinder}, \ref{chap:magnum} \\
Quotas et capacité & Chaque service plafonne les \glspl{quota} par projet ; Nova s'appuie sur Placement pour connaître la capacité réelle. & Chap.~\ref{chap:nova}, \ref{chap:placement} \\
Pilotes et greffons & Le cœur est stable, le variable est dans des pilotes : ML2 (OVS, OVN), backends de Cinder, magasins de Glance, types matériels d'Ironic, pilotes de cluster de Magnum, datastores de Trove. & Chap.~\ref{chap:neutron}, \ref{chap:cinder}, \ref{chap:ironic}, \ref{chap:magnum}, \ref{chap:trove} \\
Versions d'API négociées & Les \emph{microversions} font évoluer l'API sans casser les clients : le client indique la version qu'il attend (Nova, Placement, Cinder, Ironic). & Chap.~\ref{chap:nova}, \ref{chap:placement} \\
\bottomrule
\end{tabular}
\caption{Huit principes communs aux services d'OpenStack.}
\label{tab:concl-principes}
\end{table}

% ---------------------------------------------------------------------
\subsection*{Diagnostiquer : une méthode commune}
% ---------------------------------------------------------------------

Les principes précédents donnent aussi une méthode de diagnostic qui vaut pour tous les services (figure~\ref{fig:concl-diag}). On
part de la ressource en échec, on remonte au service qui la possède, puis on suit la requête dans ses journaux avant d'élargir aux
dépendances et à l'infrastructure.

\begin{figure}[!ht]
\centering
\adjustbox{max width=\linewidth}{%
\begin{tikzpicture}[>=Stealth, font=\footnotesize,
    cmd/.style={font=\scriptsize, text width=2.9cm, align=center}]
  \foreach \x/\t/\c [count=\i] in {1.6/{Lire\\ l'état}/{\texttt{openstack \dots{} show} : \texttt{status}, champ \texttt{fault}},
        4.9/{Trouver le\\ service}/{à qui appartient la ressource ? (\texttt{endpoint list})},
        8.2/{Suivre la\\ requête}/{identifiant \texttt{req-\dots} dans les journaux (\texttt{--debug})},
        11.5/{Vérifier les\\ dépendances}/{quotas, services et agents \texttt{up}, Placement, Glance, Neutron, Cinder},
        14.8/{Vérifier\\ l'infrastructure}/{RabbitMQ, base de données, horloge (NTP), espace disque}}
    { \node[pfoc, minimum width=2.9cm, minimum height=1.1cm] (s\i) at (\x,1.6) {\t};
      \node[num] at (\x-1.25,2.25) {\i};
      \node[cmd] at (\x,0.35) {\c}; }
  \foreach \a/\b in {1/2,2/3,3/4,4/5} \draw[fluxr] (s\a.east) -- (s\b.west);
  \node[callout, minimum width=15.6cm] at (8.2,-1.0)
    {On élargit le champ pas à pas : la cause est presque toujours dans le service propriétaire ou dans l'une de ses dépendances.};
\end{tikzpicture}}
\caption{Cinq temps pour diagnostiquer une ressource en échec, quel que soit le service.}
\label{fig:concl-diag}
\end{figure}

% ---------------------------------------------------------------------
\subsection*{Quel service pour quel besoin}
% ---------------------------------------------------------------------

Le tableau~\ref{tab:concl-besoins} part du besoin, ce qui est la démarche d'un utilisateur ou d'un architecte.

\begin{table}[!ht]
\centering
\renewcommand{\arraystretch}{1.25}
\small
\begin{tabular}{@{}>{\raggedright\arraybackslash}p{4.5cm}>{\raggedright\arraybackslash}p{3.7cm}>{\raggedright\arraybackslash}p{7.2cm}@{}}
\toprule
\textbf{Besoin} & \textbf{Service} & \textbf{À retenir} \\
\midrule
Une machine virtuelle & Nova (avec Glance, Neutron, Placement) & Disque racine local, sauf démarrage sur un volume. \\
Un disque qui survit à la machine & Cinder & Volume attaché à une instance, avec instantanés et sauvegardes. \\
Des fichiers, des archives, via HTTP & Swift & Indépendant des instances ; réplication ou codage d'effacement. \\
Réseau, routeur, adresse flottante & Neutron & Réseaux de projet isolés, groupes de sécurité. \\
Répartir le trafic entre serveurs & Octavia & Répartiteur fourni comme service, rôle des amphores. \\
Recréer une infrastructure à l'identique & Heat & Un modèle HOT décrit une pile entière. \\
Une machine physique & Ironic & Livraison de serveurs nus, souvent derrière Nova. \\
Un cluster Kubernetes & Magnum & Modèle de cluster, pilote Cluster API. \\
Une base de données managée & Trove & Instance, volume et moteur dans un conteneur. \\
Une interface web & Horizon & Appelle les API avec les droits de l'utilisateur. \\
\bottomrule
\end{tabular}
\caption{Du besoin au service.}
\label{tab:concl-besoins}
\end{table}

% ---------------------------------------------------------------------
\subsection*{Mettre en production : une liste de contrôle}
% ---------------------------------------------------------------------

Les chapitres ont chacun leur section d'exploitation ; en voici l'essentiel, dans l'ordre où l'on s'y heurte (chapitre~\ref{chap:deploiement}) :

\begin{itemize}
  \item \textbf{Topologie} : au moins trois contrôleurs, adresse virtuelle devant les API, réseaux de gestion, de tunnel et externe
        séparés.
  \item \textbf{Sécurité} : TLS sur les API exposées, secrets hors des dépôts et renouvelés, réseau de gestion interne.
  \item \textbf{Capacité} : \glspl{quota} par projet, gabarits (\emph{flavors}) et images de référence, surveillance de l'occupation des
        hyperviseurs.
  \item \textbf{Sauvegardes} : base de données de chaque service, configuration de l'outil de déploiement, volumes
        (chapitre~\ref{chap:cinder}) ; une restauration répétée au moins une fois.
  \item \textbf{Supervision} : état des services et des agents, journaux centralisés, alertes.
  \item \textbf{Pannes} : une procédure pour la perte d'un hyperviseur (évacuation) et pour la perte d'un contrôleur.
  \item \textbf{Mises à niveau} : un plan par version (versions \gls{slurp}), un environnement d'essai, une sauvegarde avant de
        commencer.
\end{itemize}

% ---------------------------------------------------------------------
\subsection*{Limites, alternatives et évolutions}
% ---------------------------------------------------------------------

\textbf{Limites.} OpenStack est puissant, mais il se paie en exploitation : de nombreux processus à superviser, un cycle de mise à
niveau tous les six mois, une documentation répartie projet par projet, une courbe d'apprentissage réelle. Cet effort se justifie
quand l'organisation en retire un bénéfice net (voir le tableau~\ref{tab:concl-alternatives}) ; il ne se justifie pas pour quelques
machines.

\begin{table}[!ht]
\centering
\renewcommand{\arraystretch}{1.25}
\small
\begin{tabular}{@{}>{\raggedright\arraybackslash}p{2.4cm}>{\raggedright\arraybackslash}p{4.4cm}>{\raggedright\arraybackslash}p{4.4cm}>{\raggedright\arraybackslash}p{4.1cm}@{}}
\toprule
\textbf{Critère} & \textbf{OpenStack (cloud privé)} & \textbf{Cloud public} & \textbf{Kubernetes seul} \\
\midrule
Ce qui est fourni & Machines, réseaux, volumes, objets, quelques services managés & Idem, avec un catalogue de services managés bien plus large & Conteneurs et leur réseau, sur des machines déjà disponibles \\
Maîtrise & Totale : matériel, données, versions & Encadrée par le contrat du fournisseur & Celle de l'infrastructure sous-jacente \\
Exploitation & À votre charge (chapitre~\ref{chap:deploiement}) & Assurée par le fournisseur & À votre charge, plus légère \\
Élasticité & Bornée par le matériel acheté & Quasi illimitée, à la demande & Bornée par les nœuds \\
Plutôt si\dots & Charge stable, souveraineté, machines physiques, équipe d'exploitation & Charge variable, démarrage rapide, peu d'équipe infrastructure & Applications déjà conteneurisées \\
\bottomrule
\end{tabular}
\caption{Trois façons de se procurer des ressources d'infrastructure (tendances générales, non des règles).}
\label{tab:concl-alternatives}
\end{table}

\textbf{Alternatives.} OpenStack et Kubernetes ne s'excluent pas, ils se complètent : Kubernetes orchestre des conteneurs sur des
machines, OpenStack fournit ces machines, leurs réseaux et leurs disques. Les deux sens de la combinaison existent. Magnum
(chapitre~\ref{chap:magnum}) fournit des clusters Kubernetes \emph{sur} OpenStack ; OpenStack-Helm, Sunbeam et RHOSO déploient OpenStack
\emph{sur} Kubernetes (section~\ref{sec:dep-autres}).

\textbf{Évolutions.} Quelques mouvements sont visibles dans le cours même :
\begin{itemize}
  \item un rythme de publication semestriel, avec des versions \gls{slurp} pour ceux qui ne mettent à niveau qu'une fois par an
        \parencite{openstack-releases,openstack-slurp} ;
  \item le départ de certains outils : TripleO est retiré, Swift n'est plus déployé par Kolla-Ansible (section~\ref{sec:dep-kolla}) ;
  \item le rapprochement avec Kubernetes, par le déploiement (Helm, Sunbeam, RHOSO) et par les services (pilotes Cluster API de Magnum,
        opérateurs de bases de données comme alternative à Trove) ;
  \item un cadre institutionnel nouveau : l'OpenInfra Foundation a rejoint la Linux Foundation le 3~juin 2025, tandis que
        OpenStack garde sa gouvernance communautaire \parencite{openinfra-lf}.
\end{itemize}

% ---------------------------------------------------------------------
\subsection*{Ce que le cours n'a pas couvert}
% ---------------------------------------------------------------------

OpenStack compte bien plus de projets officiels que les quinze chapitres \parencite{governance-projects}. Le tableau~\ref{tab:concl-autres}
en cite quelques-uns, qui complètent naturellement les services vus.

\begin{table}[!ht]
\centering
\renewcommand{\arraystretch}{1.2}
\small
\begin{tabular}{@{}>{\raggedright\arraybackslash}p{2.5cm}>{\raggedright\arraybackslash}p{5.8cm}>{\raggedright\arraybackslash}p{7.1cm}@{}}
\toprule
\textbf{Projet} & \textbf{Rôle (selon le registre officiel)} & \textbf{Lien avec le cours} \\
\midrule
Barbican & Gestionnaire de clés & Certificats des clusters de Magnum \\
Designate & Service DNS & Noms de domaine pour les instances et les adresses flottantes \\
Manila & Systèmes de fichiers partagés & Complète Cinder (blocs) et Swift (objets) \\
Skyline & Tableau de bord & Alternative à Horizon \\
Telemetry & Service de télémétrie & Mesures de l'usage des ressources, base de la supervision et de la facturation \\
Masakari & Haute disponibilité des instances & Relance des instances d'un hôte tombé (complète l'évacuation manuelle) \\
Watcher & Optimisation de l'infrastructure & Répartition de la charge entre hyperviseurs \\
Blazar & Réservation de ressources & Réserver des machines pour une période \\
Cyborg & Gestion du cycle de vie des accélérateurs & Accélérateurs matériels (GPU, FPGA) mis à disposition des instances \\
Tacker & Orchestration NFV & Fonctions réseau virtualisées, notamment chez les opérateurs de télécoms \\
CloudKitty & Service de tarification (\emph{rating}) & Facturation à partir des mesures \\
Zun & Service de conteneurs & Lancer des conteneurs sans passer par un cluster Kubernetes \\
\bottomrule
\end{tabular}
\caption{Quelques projets officiels non étudiés. Les rôles reprennent les descriptions du registre ; la colonne de droite est une lecture
d'enseignement, non une affirmation du registre.}
\label{tab:concl-autres}
\end{table}

% ---------------------------------------------------------------------
\subsection*{Pour aller plus loin}
% ---------------------------------------------------------------------

\begin{itemize}
  \item \textbf{Pratiquer} : installer DevStack ou Sunbeam sur une machine jetable (chapitre~\ref{chap:deploiement}), puis refaire le
        scénario de la figure~\ref{fig:concl-seq}, écrire un modèle Heat, ouvrir un répartiteur Octavia. Lire un journal de service
        vaut tous les schémas.
  \item \textbf{Lire} : la documentation officielle de chaque projet, le guide d'installation et le guide de haute disponibilité
        \parencite{install-guide,ha-guide}.
  \item \textbf{Se certifier} : l'examen COA (\emph{Certified OpenStack Administrator}) dure trois heures et se passe en ligne,
        sur des tâches réalisées en ligne de commande et dans Horizon ; sa page indique qu'il repose sur la version Caracal, plus
        ancienne que celles du cours, et conseille six mois d'expérience préalable \parencite{coa}.
  \item \textbf{Contribuer} : le guide du contributeur couvre le code, la documentation, les traductions et les retours
        d'opérateurs \parencite{openstack-contributors}.
\end{itemize}

% ---------------------------------------------------------------------
\subsection*{Messages clés}
% ---------------------------------------------------------------------

\begin{aretenir}
\begin{itemize}
  \item OpenStack est une famille de services indépendants, qui partagent une identité (Keystone), une infrastructure (base,
        messagerie, cache) et des conventions communes.
  \item Une action de l'utilisateur traverse plusieurs services par leurs API REST ; elle est asynchrone, et ses états disent où
        elle en est.
  \item Chaque service suit le même plan : API, coordination, agents, base, pilotes ; comprendre l'un aide à comprendre les autres.
  \item Le choix de l'outil de déploiement structure l'exploitation : Kolla-Ansible ou OpenStack-Ansible en production, DevStack pour
        apprendre et développer.
  \item Un cloud en production demande de l'exploitation : haute disponibilité, TLS, sauvegardes, supervision, mises à niveau.
  \item OpenStack complète le cloud public et Kubernetes plus qu'il ne les remplace : le bon choix dépend de la charge, de la
        maîtrise voulue et de l'équipe disponible.
\end{itemize}
\end{aretenir}

Reste à pratiquer : un cloud se comprend d'abord en le faisant tourner.
.
%  Section non numérotée, ajoutée au sommaire.
%  Sources : les chapitres du cours, le registre des projets officiels
%  (gouvernance d'OpenStack), l'annonce de l'arrivée de l'OpenInfra Foundation
%  à la Linux Foundation, la page de l'examen COA et le guide du contributeur
%  (voir sources.bib). Conclusion détaillée : 10 pages au plus.
% =====================================================================

\phantomsection
\section*{Conclusion}
\addcontentsline{toc}{section}{Conclusion}

% Figures et tableaux de cette section non numérotée : numérotation propre (C.1, C.2\dots)
% pour ne pas prolonger celle du chapitre précédent.
\setcounter{figure}{0}\setcounter{table}{0}
\renewcommand{\thefigure}{C.\arabic{figure}}\renewcommand{\thetable}{C.\arabic{table}}
\renewcommand{\theHfigure}{C.\arabic{figure}}\renewcommand{\theHtable}{C.\arabic{table}}

Ce cours a présenté OpenStack comme une famille de services indépendants qui, ensemble, transforment des serveurs, des disques et
des réseaux en un cloud en libre-service. Il est parti de la carte d'ensemble et de l'identité (Keystone, partie~\ref{part:fondations}),
a suivi la vie d'une machine virtuelle à travers le socle IaaS (Nova, Placement, Glance, Neutron, Cinder, Swift,
partie~\ref{part:socle}), puis a présenté les services qui permettent d'exploiter le cloud (Horizon, Octavia, Heat,
partie~\ref{part:exploiter}) et ceux qui fournissent d'autres ressources au-dessus du socle (Ironic, Magnum, Trove,
partie~\ref{part:plateformes}). Il s'est achevé du côté de l'administrateur, avec le déploiement (partie~\ref{part:miseenoeuvre}). Cette conclusion rassemble l'ensemble : une carte,
un scénario de bout en bout, les principes qui reviennent d'un service à l'autre, des repères pour choisir et exploiter, puis les
limites de la plateforme, ce que le cours n'a pas traité et la manière de poursuivre.

% ---------------------------------------------------------------------
\subsection*{La carte d'ensemble}
% ---------------------------------------------------------------------

La figure~\ref{fig:concl-carte} reprend la carte du chapitre~\ref{chap:presentation} (figure~\ref{fig:pres-archi}) en y ajoutant les
numéros de chapitre. Tous les services reposent sur la même infrastructure (base de données, file de
messages, cache) et sur Keystone, qui authentifie chaque appel ; l'outil de déploiement installe et met à jour le tout. Le
tableau~\ref{tab:concl-services} rappelle, pour chaque service, son rôle et les processus à connaître.

\begin{figure}[!ht]
\centering
\adjustbox{max width=\linewidth}{%
\begin{tikzpicture}[>=Stealth, font=\footnotesize]
  \node[pcli, minimum width=13.0cm] at (8.5,8.3)
    {Utilisateurs \ $\cdot$ \ ligne de commande \ $\cdot$ \ API REST \ $\cdot$ \ Horizon (chap.~\ref{chap:horizon})};
  % services d'exploitation et de plateforme
  \node[psvc, minimum width=2.6cm, minimum height=1.1cm, inner xsep=2pt] at (2.7,6.8) {\textbf{Heat} {\scriptsize\mdseries(chap.~\ref{chap:heat})}\\ {\scriptsize\mdseries orchestration}};
  \node[psvc, minimum width=2.6cm, minimum height=1.1cm, inner xsep=2pt] at (5.6,6.8) {\textbf{Octavia} {\scriptsize\mdseries(chap.~\ref{chap:octavia})}\\ {\scriptsize\mdseries répartition de charge}};
  \node[psvc, minimum width=2.6cm, minimum height=1.1cm, inner xsep=2pt] at (8.5,6.8) {\textbf{Ironic} {\scriptsize\mdseries(chap.~\ref{chap:ironic})}\\ {\scriptsize\mdseries machines physiques}};
  \node[psvc, minimum width=2.6cm, minimum height=1.1cm, inner xsep=2pt] at (11.4,6.8) {\textbf{Magnum} {\scriptsize\mdseries(chap.~\ref{chap:magnum})}\\ {\scriptsize\mdseries Kubernetes managé}};
  \node[psvc, minimum width=2.6cm, minimum height=1.1cm, inner xsep=2pt] at (14.3,6.8) {\textbf{Trove} {\scriptsize\mdseries(chap.~\ref{chap:trove})}\\ {\scriptsize\mdseries bases de données}};
  % services de base
  \node[pfoc, minimum width=3.75cm, minimum height=1.1cm] at (4.5,5.2) {Nova, Placement\\ {\scriptsize\mdseries calcul (chap.~\ref{chap:nova}, \ref{chap:placement})}};
  \node[pfoc, minimum width=3.75cm, minimum height=1.1cm] at (8.5,5.2) {Neutron\\ {\scriptsize\mdseries réseau (chap.~\ref{chap:neutron})}};
  \node[pfoc, minimum width=3.75cm, minimum height=1.1cm] at (12.5,5.2) {Glance\\ {\scriptsize\mdseries images (chap.~\ref{chap:glance})}};
  \node[psvc, minimum width=3.75cm, minimum height=1.1cm] at (6.5,3.75) {Cinder\\ {\scriptsize\mdseries volumes (chap.~\ref{chap:cinder})}};
  \node[psvc, minimum width=3.75cm, minimum height=1.1cm] at (10.5,3.75) {Swift\\ {\scriptsize\mdseries objets (chap.~\ref{chap:swift})}};
  % infrastructure
  \node[pgris, minimum width=13.0cm, minimum height=0.95cm] at (8.5,2.3)
    {Infrastructure partagée\\ {\scriptsize\mdseries MariaDB \ $\cdot$ \ RabbitMQ \ $\cdot$ \ memcached \ $\cdot$ \ HAProxy}};
  \node[pgris, minimum width=13.0cm, minimum height=0.95cm] at (8.5,1.05)
    {Machines et réseaux\\ {\scriptsize\mdseries hyperviseurs, disques, réseau physique, serveurs nus}};
  % barres latérales
  \node[draw=insarouge, thick, rounded corners=4pt, fill=insarouge!8, minimum width=1.5cm, minimum height=8.7cm] at (0.45,4.65) {};
  \node[rotate=90, font=\footnotesize\bfseries, text=insarouge] at (0.45,4.65) {Keystone : identité, catalogue, tokens (chap.~\ref{chap:keystone})};
  \node[draw=insarouge, thick, rounded corners=4pt, fill=insarouge!8, minimum width=1.5cm, minimum height=8.7cm] at (16.55,4.65) {};
  \node[rotate=90, font=\footnotesize\bfseries, text=insarouge] at (16.55,4.65) {Déploiement et exploitation (chap.~\ref{chap:deploiement})};
  \node[callout, minimum width=16.4cm] at (8.5,-0.55)
    {Des services indépendants, une identité commune, une infrastructure partagée : la plateforme tient à ces trois choix.};
\end{tikzpicture}}
\caption{Les services d'OpenStack étudiés dans le cours. En haut, les services d'exploitation et de plateforme (parties~\ref{part:exploiter}
et~\ref{part:plateformes}) ; au milieu, le socle IaaS (partie~\ref{part:socle}), en rouge plein pour le cœur d'un cloud (calcul, réseau,
images) ; en gris, ce qui les supporte.}
\label{fig:concl-carte}
\end{figure}

\begin{table}[!ht]
\centering
\renewcommand{\arraystretch}{1.25}
\small
\begin{tabular}{@{}>{\raggedright\arraybackslash}p{2.0cm}>{\raggedright\arraybackslash}p{3.9cm}>{\raggedright\arraybackslash}p{8.3cm}r@{}}
\toprule
\textbf{Service} & \textbf{Rôle} & \textbf{Processus et composants à connaître} & \textbf{Chap.} \\
\midrule
Keystone & Identité, authentification, catalogue & Service HTTP (WSGI) ; tokens Fernet ; catalogue des points d'accès & \ref{chap:keystone} \\
Nova & Instances (calcul) & \texttt{nova-api}, \texttt{nova-conductor}, \texttt{nova-scheduler}, \texttt{nova-compute} & \ref{chap:nova} \\
Placement & Inventaire et allocation des ressources & \texttt{placement-api} : fournisseurs, classes, allocations & \ref{chap:placement} \\
Glance & Images & \texttt{glance-api} et ses magasins & \ref{chap:glance} \\
Neutron & Réseau & \texttt{neutron-server}, agents (L3, DHCP), OVS ou OVN & \ref{chap:neutron} \\
Cinder & Volumes (stockage bloc) & \texttt{cinder-api}, \texttt{-scheduler}, \texttt{-volume}, \texttt{-backup} & \ref{chap:cinder} \\
Swift & Stockage objet & Proxy, serveurs de comptes, de conteneurs et d'objets ; anneaux & \ref{chap:swift} \\
Horizon & Tableau de bord & Application Django sous Apache (\texttt{mod\_wsgi}) & \ref{chap:horizon} \\
Octavia & Répartition de charge & \texttt{octavia-api}, worker, health manager, housekeeping ; amphores & \ref{chap:octavia} \\
Heat & Orchestration & \texttt{heat-api}, \texttt{heat-engine} ; modèles HOT & \ref{chap:heat} \\
Ironic & Machines physiques & \texttt{ironic-api}, \texttt{ironic-conductor}, agent \texttt{ironic-python-agent} & \ref{chap:ironic} \\
Magnum & Clusters Kubernetes & \texttt{magnum-api}, \texttt{magnum-conductor} ; pilote de cluster & \ref{chap:magnum} \\
Trove & Bases de données & \texttt{trove-api}, \texttt{trove-taskmanager}, \texttt{trove-conductor} ; agent invité & \ref{chap:trove} \\
Outils & Installer, exploiter, mettre à niveau & Kolla-Ansible, OpenStack-Ansible, DevStack, \dots & \ref{chap:deploiement} \\
\bottomrule
\end{tabular}
\caption{Les services du cours en un coup d'œil.}
\label{tab:concl-services}
\end{table}

% ---------------------------------------------------------------------
\subsection*{Un scénario de bout en bout}
% ---------------------------------------------------------------------

Rien ne montre mieux la coopération des services qu'une demande ordinaire : une machine virtuelle dotée d'un disque de données et
d'une adresse publique. Le chapitre~\ref{chap:presentation} avait suivi la création d'une machine seule (figure~\ref{fig:pres-scenario}) ; ici, cinq
commandes y ajoutent un volume et une adresse publique, et la figure~\ref{fig:concl-seq} en suit les appels.

\par\noindent\begin{minipage}{\linewidth}
\begin{lstlisting}[style=terminal]
$ openstack volume create --size 20 donnees
$ openstack server create --image ubuntu-24.04 --flavor m1.small --network prive vm1
$ openstack server add volume vm1 donnees
$ openstack floating ip create public
$ openstack server add floating ip vm1 203.0.113.10
\end{lstlisting}
\end{minipage}\par

\begin{figure}[!ht]
\centering
\adjustbox{max width=\linewidth}{%
\begin{tikzpicture}[>=Stealth, font=\footnotesize]
  \node[acteur, minimum width=2.4cm] (c) at (0.8,0)  {Client};
  \node[acteur, minimum width=2.4cm, fill=insapeche, draw=insaorange] (k) at (3.9,0)  {Keystone};
  \node[acteur, minimum width=2.4cm, fill=insapeche, draw=insaorange] (ci) at (7.0,0)  {Cinder};
  \node[acteur, minimum width=2.4cm, fill=insarouge, text=white, draw=black] (n) at (10.1,0) {Nova};
  \node[acteur, minimum width=2.4cm, fill=insapeche, draw=insaorange] (g) at (13.2,0) {Glance};
  \node[acteur, minimum width=2.4cm, fill=insapeche, draw=insaorange] (ne) at (16.3,0) {Neutron};
  \foreach \n in {c,k,ci,n,g,ne} \draw[vie] (\n.south) -- ++(0,-10.9);
  \seqmsg{0.8}{3.9}{-1.3}{1}{authentification}
  \seqrep{3.9}{0.8}{-2.0}{}{token}
  \seqmsg{0.8}{7.0}{-2.9}{2}{\texttt{POST /volumes}}
  \seqrep{7.0}{0.8}{-3.6}{}{volume \texttt{available}}
  \seqmsg{0.8}{10.1}{-4.5}{3}{\texttt{POST /servers}}
  \seqrep{10.1}{0.8}{-5.2}{}{\texttt{202 Accepted}}
  \seqmsg{10.1}{13.2}{-6.1}{4}{image}
  \seqmsg{10.1}{16.3}{-7.0}{5}{port réseau}
  \node[callout, font=\scriptsize] at (10.1,-7.8) {instance \texttt{ACTIVE}};
  \seqmsg{0.8}{10.1}{-8.6}{6}{\texttt{add volume}}
  \seqmsg{10.1}{7.0}{-9.5}{7}{attacher le volume}
  \seqmsg{0.8}{16.3}{-10.4}{8}{adresse flottante}
  \node[callout, font=\scriptsize, text width=4.4cm] at (14.6,-2.9)
    {Chaque appel porte le token ; chaque service vérifie les droits et les quotas.};
\end{tikzpicture}}
\caption{Le scénario en huit échanges. Le client s'authentifie (1), crée le volume (2) et la machine (3) ; Nova répond sans attendre,
puis obtient l'image et le port réseau auprès de Glance et de Neutron (4, 5), détaillés dans la figure~\ref{fig:nova-seq}. Le
volume est attaché par Nova auprès de Cinder (6, 7) et l'adresse flottante associée au port par Neutron (8).}
\label{fig:concl-seq}
\end{figure}

Quatre idées se dégagent. Un seul token, émis par Keystone, sert auprès de tous les services. Les services collaborent par leurs API
REST, sans jamais lire la base les uns des autres. Nova répond \texttt{202 Accepted} avant d'avoir terminé : l'état de l'instance
(\texttt{BUILD}, puis \texttt{ACTIVE}) renseigne sur l'avancement. Enfin, quand une machine finit en \texttt{ERROR}, la bonne question est
« quel service a refusé, et pourquoi ? » (quota, image, port, volume) : on lit l'état de la ressource, puis les journaux du service
concerné.

% ---------------------------------------------------------------------
\subsection*{Les principes qui reviennent partout}
% ---------------------------------------------------------------------

Au-delà de leurs différences, les services partagent une même manière de fonctionner (tableau~\ref{tab:concl-principes}). La
reconnaître dans un service inconnu permet de le comprendre vite : on cherche son API, son processus de coordination, ses agents,
sa base, ses états et ses pilotes. Les bibliothèques \gls{oslo} (configuration, messagerie, accès aux bases, politiques) en sont
le ciment commun.

\begin{table}[!ht]
\centering
\renewcommand{\arraystretch}{1.25}
\small
\begin{tabular}{@{}>{\raggedright\arraybackslash}p{3.2cm}>{\raggedright\arraybackslash}p{9.0cm}>{\raggedright\arraybackslash}p{3.1cm}@{}}
\toprule
\textbf{Principe} & \textbf{En pratique} & \textbf{À revoir} \\
\midrule
API REST authentifiée & Chaque appel porte un token de Keystone ; le service le valide et applique sa politique de droits. & Chap.~\ref{chap:keystone} \\
API, coordination, agents & L'API valide et répond ; un processus de coordination (\emph{conductor}, \emph{engine}, \emph{taskmanager}) décide ; des agents exécutent sur les nœuds. & Chap.~\ref{chap:nova}, \ref{chap:cinder}, \ref{chap:heat}, \ref{chap:ironic}, \ref{chap:trove} \\
Messagerie interne & Les processus d'un même service s'appellent par RPC sur la file de messages (RabbitMQ) ; d'un service à l'autre, on passe par l'API REST. & Chap.~\ref{chap:nova}, \ref{chap:deploiement} \\
Une base par service & Chaque service gère son schéma (\texttt{db sync}) ; les autres ne le lisent jamais. & Chap.~\ref{chap:deploiement} \\
Opérations asynchrones & L'API répond vite ; la ressource passe par des états (\texttt{BUILD} $\to$ \texttt{ACTIVE}, \texttt{creating} $\to$ \texttt{available}) et finit en \texttt{ERROR} s'il y a un échec. & Chap.~\ref{chap:nova}, \ref{chap:cinder}, \ref{chap:magnum} \\
Quotas et capacité & Chaque service plafonne les \glspl{quota} par projet ; Nova s'appuie sur Placement pour connaître la capacité réelle. & Chap.~\ref{chap:nova}, \ref{chap:placement} \\
Pilotes et greffons & Le cœur est stable, le variable est dans des pilotes : ML2 (OVS, OVN), backends de Cinder, magasins de Glance, types matériels d'Ironic, pilotes de cluster de Magnum, datastores de Trove. & Chap.~\ref{chap:neutron}, \ref{chap:cinder}, \ref{chap:ironic}, \ref{chap:magnum}, \ref{chap:trove} \\
Versions d'API négociées & Les \emph{microversions} font évoluer l'API sans casser les clients : le client indique la version qu'il attend (Nova, Placement, Cinder, Ironic). & Chap.~\ref{chap:nova}, \ref{chap:placement} \\
\bottomrule
\end{tabular}
\caption{Huit principes communs aux services d'OpenStack.}
\label{tab:concl-principes}
\end{table}

% ---------------------------------------------------------------------
\subsection*{Diagnostiquer : une méthode commune}
% ---------------------------------------------------------------------

Les principes précédents donnent aussi une méthode de diagnostic qui vaut pour tous les services (figure~\ref{fig:concl-diag}). On
part de la ressource en échec, on remonte au service qui la possède, puis on suit la requête dans ses journaux avant d'élargir aux
dépendances et à l'infrastructure.

\begin{figure}[!ht]
\centering
\adjustbox{max width=\linewidth}{%
\begin{tikzpicture}[>=Stealth, font=\footnotesize,
    cmd/.style={font=\scriptsize, text width=2.9cm, align=center}]
  \foreach \x/\t/\c [count=\i] in {1.6/{Lire\\ l'état}/{\texttt{openstack \dots{} show} : \texttt{status}, champ \texttt{fault}},
        4.9/{Trouver le\\ service}/{à qui appartient la ressource ? (\texttt{endpoint list})},
        8.2/{Suivre la\\ requête}/{identifiant \texttt{req-\dots} dans les journaux (\texttt{--debug})},
        11.5/{Vérifier les\\ dépendances}/{quotas, services et agents \texttt{up}, Placement, Glance, Neutron, Cinder},
        14.8/{Vérifier\\ l'infrastructure}/{RabbitMQ, base de données, horloge (NTP), espace disque}}
    { \node[pfoc, minimum width=2.9cm, minimum height=1.1cm] (s\i) at (\x,1.6) {\t};
      \node[num] at (\x-1.25,2.25) {\i};
      \node[cmd] at (\x,0.35) {\c}; }
  \foreach \a/\b in {1/2,2/3,3/4,4/5} \draw[fluxr] (s\a.east) -- (s\b.west);
  \node[callout, minimum width=15.6cm] at (8.2,-1.0)
    {On élargit le champ pas à pas : la cause est presque toujours dans le service propriétaire ou dans l'une de ses dépendances.};
\end{tikzpicture}}
\caption{Cinq temps pour diagnostiquer une ressource en échec, quel que soit le service.}
\label{fig:concl-diag}
\end{figure}

% ---------------------------------------------------------------------
\subsection*{Quel service pour quel besoin}
% ---------------------------------------------------------------------

Le tableau~\ref{tab:concl-besoins} part du besoin, ce qui est la démarche d'un utilisateur ou d'un architecte.

\begin{table}[!ht]
\centering
\renewcommand{\arraystretch}{1.25}
\small
\begin{tabular}{@{}>{\raggedright\arraybackslash}p{4.5cm}>{\raggedright\arraybackslash}p{3.7cm}>{\raggedright\arraybackslash}p{7.2cm}@{}}
\toprule
\textbf{Besoin} & \textbf{Service} & \textbf{À retenir} \\
\midrule
Une machine virtuelle & Nova (avec Glance, Neutron, Placement) & Disque racine local, sauf démarrage sur un volume. \\
Un disque qui survit à la machine & Cinder & Volume attaché à une instance, avec instantanés et sauvegardes. \\
Des fichiers, des archives, via HTTP & Swift & Indépendant des instances ; réplication ou codage d'effacement. \\
Réseau, routeur, adresse flottante & Neutron & Réseaux de projet isolés, groupes de sécurité. \\
Répartir le trafic entre serveurs & Octavia & Répartiteur fourni comme service, rôle des amphores. \\
Recréer une infrastructure à l'identique & Heat & Un modèle HOT décrit une pile entière. \\
Une machine physique & Ironic & Livraison de serveurs nus, souvent derrière Nova. \\
Un cluster Kubernetes & Magnum & Modèle de cluster, pilote Cluster API. \\
Une base de données managée & Trove & Instance, volume et moteur dans un conteneur. \\
Une interface web & Horizon & Appelle les API avec les droits de l'utilisateur. \\
\bottomrule
\end{tabular}
\caption{Du besoin au service.}
\label{tab:concl-besoins}
\end{table}

% ---------------------------------------------------------------------
\subsection*{Mettre en production : une liste de contrôle}
% ---------------------------------------------------------------------

Les chapitres ont chacun leur section d'exploitation ; en voici l'essentiel, dans l'ordre où l'on s'y heurte (chapitre~\ref{chap:deploiement}) :

\begin{itemize}
  \item \textbf{Topologie} : au moins trois contrôleurs, adresse virtuelle devant les API, réseaux de gestion, de tunnel et externe
        séparés.
  \item \textbf{Sécurité} : TLS sur les API exposées, secrets hors des dépôts et renouvelés, réseau de gestion interne.
  \item \textbf{Capacité} : \glspl{quota} par projet, gabarits (\emph{flavors}) et images de référence, surveillance de l'occupation des
        hyperviseurs.
  \item \textbf{Sauvegardes} : base de données de chaque service, configuration de l'outil de déploiement, volumes
        (chapitre~\ref{chap:cinder}) ; une restauration répétée au moins une fois.
  \item \textbf{Supervision} : état des services et des agents, journaux centralisés, alertes.
  \item \textbf{Pannes} : une procédure pour la perte d'un hyperviseur (évacuation) et pour la perte d'un contrôleur.
  \item \textbf{Mises à niveau} : un plan par version (versions \gls{slurp}), un environnement d'essai, une sauvegarde avant de
        commencer.
\end{itemize}

% ---------------------------------------------------------------------
\subsection*{Limites, alternatives et évolutions}
% ---------------------------------------------------------------------

\textbf{Limites.} OpenStack est puissant, mais il se paie en exploitation : de nombreux processus à superviser, un cycle de mise à
niveau tous les six mois, une documentation répartie projet par projet, une courbe d'apprentissage réelle. Cet effort se justifie
quand l'organisation en retire un bénéfice net (voir le tableau~\ref{tab:concl-alternatives}) ; il ne se justifie pas pour quelques
machines.

\begin{table}[!ht]
\centering
\renewcommand{\arraystretch}{1.25}
\small
\begin{tabular}{@{}>{\raggedright\arraybackslash}p{2.4cm}>{\raggedright\arraybackslash}p{4.4cm}>{\raggedright\arraybackslash}p{4.4cm}>{\raggedright\arraybackslash}p{4.1cm}@{}}
\toprule
\textbf{Critère} & \textbf{OpenStack (cloud privé)} & \textbf{Cloud public} & \textbf{Kubernetes seul} \\
\midrule
Ce qui est fourni & Machines, réseaux, volumes, objets, quelques services managés & Idem, avec un catalogue de services managés bien plus large & Conteneurs et leur réseau, sur des machines déjà disponibles \\
Maîtrise & Totale : matériel, données, versions & Encadrée par le contrat du fournisseur & Celle de l'infrastructure sous-jacente \\
Exploitation & À votre charge (chapitre~\ref{chap:deploiement}) & Assurée par le fournisseur & À votre charge, plus légère \\
Élasticité & Bornée par le matériel acheté & Quasi illimitée, à la demande & Bornée par les nœuds \\
Plutôt si\dots & Charge stable, souveraineté, machines physiques, équipe d'exploitation & Charge variable, démarrage rapide, peu d'équipe infrastructure & Applications déjà conteneurisées \\
\bottomrule
\end{tabular}
\caption{Trois façons de se procurer des ressources d'infrastructure (tendances générales, non des règles).}
\label{tab:concl-alternatives}
\end{table}

\textbf{Alternatives.} OpenStack et Kubernetes ne s'excluent pas, ils se complètent : Kubernetes orchestre des conteneurs sur des
machines, OpenStack fournit ces machines, leurs réseaux et leurs disques. Les deux sens de la combinaison existent. Magnum
(chapitre~\ref{chap:magnum}) fournit des clusters Kubernetes \emph{sur} OpenStack ; OpenStack-Helm, Sunbeam et RHOSO déploient OpenStack
\emph{sur} Kubernetes (section~\ref{sec:dep-autres}).

\textbf{Évolutions.} Quelques mouvements sont visibles dans le cours même :
\begin{itemize}
  \item un rythme de publication semestriel, avec des versions \gls{slurp} pour ceux qui ne mettent à niveau qu'une fois par an
        \parencite{openstack-releases,openstack-slurp} ;
  \item le départ de certains outils : TripleO est retiré, Swift n'est plus déployé par Kolla-Ansible (section~\ref{sec:dep-kolla}) ;
  \item le rapprochement avec Kubernetes, par le déploiement (Helm, Sunbeam, RHOSO) et par les services (pilotes Cluster API de Magnum,
        opérateurs de bases de données comme alternative à Trove) ;
  \item un cadre institutionnel nouveau : l'OpenInfra Foundation a rejoint la Linux Foundation le 3~juin 2025, tandis que
        OpenStack garde sa gouvernance communautaire \parencite{openinfra-lf}.
\end{itemize}

% ---------------------------------------------------------------------
\subsection*{Ce que le cours n'a pas couvert}
% ---------------------------------------------------------------------

OpenStack compte bien plus de projets officiels que les quinze chapitres \parencite{governance-projects}. Le tableau~\ref{tab:concl-autres}
en cite quelques-uns, qui complètent naturellement les services vus.

\begin{table}[!ht]
\centering
\renewcommand{\arraystretch}{1.2}
\small
\begin{tabular}{@{}>{\raggedright\arraybackslash}p{2.5cm}>{\raggedright\arraybackslash}p{5.8cm}>{\raggedright\arraybackslash}p{7.1cm}@{}}
\toprule
\textbf{Projet} & \textbf{Rôle (selon le registre officiel)} & \textbf{Lien avec le cours} \\
\midrule
Barbican & Gestionnaire de clés & Certificats des clusters de Magnum \\
Designate & Service DNS & Noms de domaine pour les instances et les adresses flottantes \\
Manila & Systèmes de fichiers partagés & Complète Cinder (blocs) et Swift (objets) \\
Skyline & Tableau de bord & Alternative à Horizon \\
Telemetry & Service de télémétrie & Mesures de l'usage des ressources, base de la supervision et de la facturation \\
Masakari & Haute disponibilité des instances & Relance des instances d'un hôte tombé (complète l'évacuation manuelle) \\
Watcher & Optimisation de l'infrastructure & Répartition de la charge entre hyperviseurs \\
Blazar & Réservation de ressources & Réserver des machines pour une période \\
Cyborg & Gestion du cycle de vie des accélérateurs & Accélérateurs matériels (GPU, FPGA) mis à disposition des instances \\
Tacker & Orchestration NFV & Fonctions réseau virtualisées, notamment chez les opérateurs de télécoms \\
CloudKitty & Service de tarification (\emph{rating}) & Facturation à partir des mesures \\
Zun & Service de conteneurs & Lancer des conteneurs sans passer par un cluster Kubernetes \\
\bottomrule
\end{tabular}
\caption{Quelques projets officiels non étudiés. Les rôles reprennent les descriptions du registre ; la colonne de droite est une lecture
d'enseignement, non une affirmation du registre.}
\label{tab:concl-autres}
\end{table}

% ---------------------------------------------------------------------
\subsection*{Pour aller plus loin}
% ---------------------------------------------------------------------

\begin{itemize}
  \item \textbf{Pratiquer} : installer DevStack ou Sunbeam sur une machine jetable (chapitre~\ref{chap:deploiement}), puis refaire le
        scénario de la figure~\ref{fig:concl-seq}, écrire un modèle Heat, ouvrir un répartiteur Octavia. Lire un journal de service
        vaut tous les schémas.
  \item \textbf{Lire} : la documentation officielle de chaque projet, le guide d'installation et le guide de haute disponibilité
        \parencite{install-guide,ha-guide}.
  \item \textbf{Se certifier} : l'examen COA (\emph{Certified OpenStack Administrator}) dure trois heures et se passe en ligne,
        sur des tâches réalisées en ligne de commande et dans Horizon ; sa page indique qu'il repose sur la version Caracal, plus
        ancienne que celles du cours, et conseille six mois d'expérience préalable \parencite{coa}.
  \item \textbf{Contribuer} : le guide du contributeur couvre le code, la documentation, les traductions et les retours
        d'opérateurs \parencite{openstack-contributors}.
\end{itemize}

% ---------------------------------------------------------------------
\subsection*{Messages clés}
% ---------------------------------------------------------------------

\begin{aretenir}
\begin{itemize}
  \item OpenStack est une famille de services indépendants, qui partagent une identité (Keystone), une infrastructure (base,
        messagerie, cache) et des conventions communes.
  \item Une action de l'utilisateur traverse plusieurs services par leurs API REST ; elle est asynchrone, et ses états disent où
        elle en est.
  \item Chaque service suit le même plan : API, coordination, agents, base, pilotes ; comprendre l'un aide à comprendre les autres.
  \item Le choix de l'outil de déploiement structure l'exploitation : Kolla-Ansible ou OpenStack-Ansible en production, DevStack pour
        apprendre et développer.
  \item Un cloud en production demande de l'exploitation : haute disponibilité, TLS, sauvegardes, supervision, mises à niveau.
  \item OpenStack complète le cloud public et Kubernetes plus qu'il ne les remplace : le bon choix dépend de la charge, de la
        maîtrise voulue et de l'équipe disponible.
\end{itemize}
\end{aretenir}

Reste à pratiquer : un cloud se comprend d'abord en le faisant tourner.
